% 04 — FRAGMENTO SUCESOR HODGE Y BIRCH--SWINNERTON-DYER
% =======================================================
% Inserción autónoma prevista tras el marcador de integración del Libro IV.
% No importa fuentes probatorios ni modifica las tres realizaciones
% precedentes.
%
% Trazabilidad técnica no narrativa:
%   Hodge: CHI-II-HDG-01/02/03, CHI-IV-HDG-01/C01.
%   BSD:   CHI-II-BSD-01/02/03, CHI-IV-BSD-01/02/C01.
%   Fuentes: celda APP--Mayer--Vietoris, totalizador coend y controles
%   elípticos; producto fibrado genealógico de Manin, publicación acoplada
%   Kato--PT, dualidad ortogonal reducida, relleno finita--singular y regulador
%   Bockstein derivado. Los hashes y estados de ejecución permanecen en el
%   expediente técnico de la edición sucesora.

\ifdefined\HMTSucesoraStructureTrace
  \typeout{HMT-SUCESORA 04: realizaciones Hodge y Birch--Swinnerton-Dyer}
\fi

\chapter{Completitud de la historia cohomológica enriquecida y algebraicidad de las clases racionales de Hodge}
\chaptermark{Completitud cohomológica y algebraicidad de Hodge}

\subsection*{Objeto, dominio y datos de partida}

Sea \(X\) una variedad proyectiva lisa sobre \(\mathbb C\) y sea
\(p\geq0\). Se escribe
\begin{equation}\label{eq:hodge-suc-dominio}
 \operatorname{Hdg}^{p}(X)_{\mathbb Q}
 :=H^{2p}(X,\mathbb Q)
   \cap H^{p,p}(X,\mathbb C).
\end{equation}
Sea \(\mathbf{Perf}^{\geq p}(X)\) la categoría idempotente completa de
complejos perfectos cuyo soporte tiene codimensión al menos \(p\). En esta
sección se emplea la abreviatura
\[
 \operatorname{Perf}^{\geq p}(X)_{\mathbb Q}
 :=K_0\!\left(\mathbf{Perf}^{\geq p}(X)\right)
   \otimes_{\mathbb Z}\mathbb Q;
\]
por tanto, \(R_{\mathrm{Hdg}}\) publica la clase del complejo perfecto
construido. El punto de partida no es una clase desnuda de
\eqref{eq:hodge-suc-dominio}, sino el límite de historias supervivientes
\begin{equation}\label{eq:hodge-suc-xchi}
 X_{\chi}(X,p)
 =\varprojlim_{K}\operatorname{Surv}_{K}(X,p).
\end{equation}
Cada elemento de \(X_{\chi}(X,p)\) conserva hoja, orientación, cociente,
acarreo, frontera, ruta y memoria. En particular, el retorno de la conexión
nonádica conjunta satisface
\begin{equation}\label{eq:hodge-suc-retorno}
 \Gamma _9=\operatorname{Hol}_{C_9}(\nabla_{\mathrm{TPK}}),
 \qquad g(k+9)=g(k),
 \qquad \widehat x(k+9)\neq\widehat x(k)
 \quad\text{en general}.
\end{equation}
La igualdad conserva la fase; la desigualdad conserva la profundidad de la
historia. Esta distinción excluye que un retorno se interprete como una
reinicialización.

La composición de rutas está gobernada por el cociclo de acarreo
\begin{equation}\label{eq:hodge-suc-carry}
 \operatorname{Carr}(\gamma _2\gamma _1)
 =\operatorname{Carr}(\gamma _2)H(\gamma _1)
  +Y(\gamma _2)\operatorname{Carr}(\gamma _1),
\end{equation}
y la memoria acumulada por
\begin{equation}\label{eq:hodge-suc-memoria}
 S_0=\sum_{r<N}(M_9^{*})^{r}D_9^{*}D_9M_9^{r}
     +(M_9^{*})^{N}S_0M_9^{N}.
\end{equation}
Las ecuaciones \eqref{eq:hodge-suc-xchi}--\eqref{eq:hodge-suc-memoria}
son datos construidos en la Parte II. Aquí se especializan; no se
reconstruyen a partir del codominio cohomológico.

\subsection*{Celda local APP--Mayer--Vietoris}

Sean \(Z,W\subset X\) incidencias cerradas con intersección
Tor-independiente y soportes orientados de la codimensión declarada. Para
enteros \(a,b\geq1\), defínase
\begin{equation}\label{eq:hodge-suc-kappa}
 \begin{aligned}
 \kappa_{a,b}:\mathcal O_Z^{a}\oplus\mathcal O_W^{b}
 &\longrightarrow \mathcal O_{Z\cap W}^{ab},\\
 ((u_i)_{i=1}^{a},(v_j)_{j=1}^{b})
 &\longmapsto
 \bigl(u_i|_{Z\cap W}-v_j|_{Z\cap W}\bigr)_{i,j}.
 \end{aligned}
\end{equation}
Sobre una fibra de \(Z\cap W\), el núcleo diagonal de
\(\kappa_{a,b}\) tiene rango uno y su conúcleo tiene rango
\begin{equation}\label{eq:hodge-suc-defecto-rango}
 ab-(a+b-1)=(a-1)(b-1).
\end{equation}
Si \((\rho_{\Sigma},c_{\Sigma})\) y
\((\rho_{\Pi},c_{\Pi})\) son, respectivamente, los residuos y acarreos de
las hojas aditiva y multiplicativa de APP, la misma cantidad se expresa por
\begin{equation}\label{eq:hodge-suc-celda-carry}
 (a-1)(b-1)
 =\rho_{\Pi}-\rho_{\Sigma}+1
  +9\bigl(c_{\Pi}-c_{\Sigma}\bigr).
\end{equation}
Así, el defecto local no es una corrección exterior: es la publicación en
\(\operatorname{Perf}\) de la diferencia entre ambas hojas con su acarreo.

Sea \(\tau\) el intercambio de los dos sumandos de la fuente y sea \(P\)
la transposición de los índices del blanco. La orientación TRIT fija
\begin{equation}\label{eq:hodge-suc-orientacion-kappa}
 \kappa_{b,a}\tau=-P\kappa_{a,b}.
\end{equation}
El signo de \eqref{eq:hodge-suc-orientacion-kappa} forma parte del objeto
orientado. Su omisión identificaría rutas opuestas y destruiría la
naturalidad del carácter localizado.

\begin{proposition}[Realización perfecta de la celda local]
\label{prop:hodge-suc-celda-perf}
El cono
\begin{equation}\label{eq:hodge-suc-cone-kappa}
 \mathcal C_{a,b}(Z,W):=\operatorname{Cone}(\kappa_{a,b})
\end{equation}
es un complejo perfecto con el soporte prescrito. Su clase localizada
conserva el defecto de rango \eqref{eq:hodge-suc-defecto-rango}, el acarreo
\eqref{eq:hodge-suc-celda-carry} y la orientación
\eqref{eq:hodge-suc-orientacion-kappa}.
\end{proposition}

\begin{proof}
La Tor-independencia identifica la restricción derivada al cruce con la
restricción ordinaria empleada en \eqref{eq:hodge-suc-kappa}. Las fuentes y
el blanco son perfectos sobre sus soportes; por estabilidad de
\(\operatorname{Perf}\) bajo conos, también lo es
\(\mathcal C_{a,b}(Z,W)\). El cálculo fibra a fibra da
\eqref{eq:hodge-suc-defecto-rango}; la división APP de los datos de suma y
producto da \eqref{eq:hodge-suc-celda-carry}. Finalmente, intercambiar
\(Z,a\) con \(W,b\) cambia cada diferencia por su opuesta, lo que prueba
\eqref{eq:hodge-suc-orientacion-kappa}. Como el carácter de Chern localizado
es aditivo en triángulos distinguidos, las tres informaciones pasan a la
clase del cono.
\end{proof}

\subsection*{Totalización de historias y conservación del terminal}

Sea \(J_N\) la categoría finita de historias compatibles a profundidad
\(N\), sea \(\Gamma\) el diagrama de complejos de incidencia obtenido de
\eqref{eq:hodge-suc-cone-kappa} y sea \(\mathsf W\) el módulo derecho que
registra las orientaciones y pesos de ruta. La totalización derivada es el
coend
\begin{equation}\label{eq:hodge-suc-coend}
 K_{X,p,N}
 :=\mathsf W\otimes^{\mathbf L}_{\mathbb Z[J_N]}\Gamma.
\end{equation}
Los enlaces de supervivencia forman un sistema compatible en \(N\). La
holonomía \(\Gamma_9\) induce un endomorfismo
\(U_{\Gamma_9}\) del totalizador y el objeto terminal se conserva como
\begin{equation}\label{eq:hodge-suc-terminal}
 \mathcal T_{X,p}
 :=\operatorname{Cone}\bigl(I-U_{\Gamma_9}\bigr).
\end{equation}
El cono \eqref{eq:hodge-suc-terminal} impide que la totalización confunda
una fase repetida con la misma historia: el valor visible puede retornar,
pero la diferencia terminal retiene el desplazamiento de memoria.

La compatibilidad entre \eqref{eq:hodge-suc-coend} y el límite
\eqref{eq:hodge-suc-xchi} define el morfismo de realización
\begin{equation}\label{eq:hodge-suc-realizacion}
 R_{\mathrm{Hdg}}:
 X_{\chi}(X,p)\longrightarrow
 \operatorname{Perf}^{\geq p}(X)_{\mathbb Q}.
\end{equation}
Este morfismo se fija en la fuente: no recibe como entrada una clase de
Hodge ni un representante que deba reconocer.

\begin{theorem}[Identidad de publicación Hodge]
\label{thm:hodge-suc-identidad}
Sobre \(X_{\chi}(X,p)\) conmuta la identidad
\begin{equation}\label{eq:hodge-suc-identidad}
 \operatorname{cl}_{X,p}\circ
 \operatorname{ch}^{\mathrm{loc}}_{p}\circ
 R_{\mathrm{Hdg}}
 =\operatorname{ev}^{\mathrm{Hdg}}_{\infty,X,p}.
\end{equation}
Aquí
\(\operatorname{ch}^{\mathrm{loc}}_{p}\) toma la componente de
codimensión \(p\) del carácter localizado y
\(\operatorname{cl}_{X,p}\) es su publicación cohomológica.
\end{theorem}

\begin{proof}
En una celda, la igualdad es la compatibilidad del carácter localizado con
la restricción, el triángulo de Mayer--Vietoris y la clase cohomológica. La
orientación está fijada por \eqref{eq:hodge-suc-orientacion-kappa}; por
tanto, no queda un signo por elegir. La aditividad del carácter de Chern
transporta la igualdad a cada cono \eqref{eq:hodge-suc-cone-kappa}. La
naturalidad respecto de las flechas de \(J_N\) la hace descender al coend
\eqref{eq:hodge-suc-coend}. La ley de acarreo
\eqref{eq:hodge-suc-carry} garantiza su compatibilidad bajo composición y
\eqref{eq:hodge-suc-terminal} conserva el término terminal. Al pasar al
límite proyectivo se obtienen dos transformaciones naturales con iguales
componentes en todo cilindro finito; luego coinciden y resulta
\eqref{eq:hodge-suc-identidad}.
\end{proof}

El teorema de completitud del atlas de historias enriquecidas, aplicado a
esta especialización, establece
\begin{equation}\label{eq:hodge-suc-completitud}
 \operatorname{ev}^{\mathrm{Hdg}}_{\infty,X,p}
 \bigl(X_{\chi}(X,p)\bigr)
 =\operatorname{Hdg}^{p}(X)_{\mathbb Q}.
\end{equation}
Ésta es la estructura de partida explícita de la especialización Hodge.

% Ampliación sustantiva de la completitud coinductiva Hodge.
% Módulo destinado al capítulo Hodge de la Parte IV.

\section{Despliegue coinductivo de la completitud Hodge}
\label{sec:amp-hodge-completitud-desplegada}

La igualdad de completitud
\eqref{eq:hodge-suc-completitud} puede desplegarse sin comenzar por una clase
algebraica ni introducir en la fuente la conclusión que se desea publicar. El
objeto que se completa es el sistema de historias enriquecidas; la clase de
Hodge interviene únicamente como dato de evaluación compatible en sus
truncaciones finitas.

\subsection{Sistemas finitos de observación y fibras compatibles}
\label{subsec:amp-hodge-sistemas-finitos}

Para cada profundidad \(N\geq0\), sea \(J_N(X,p)\) la categoría finita de
cartas de incidencia, rutas orientadas y terminales de memoria de profundidad a
lo sumo \(N\). Se escribe
\begin{equation}\label{eq:amp-hodge-surv-n}
 \mathscr S_N(X,p):=\operatorname{Surv}_{J_N}(X,p),
 \qquad
 \rho_{N+1,N}:\mathscr S_{N+1}(X,p)\longrightarrow\mathscr S_N(X,p)
\end{equation}
para el espacio racionalizado de estados supervivientes y su restricción. Los
lectores finitos toman valores en el módulo de observaciones de cilindro
\(\mathscr H_N^p(X)\):
\begin{equation}\label{eq:amp-hodge-evaluacion-finita}
 e_N^{\mathrm{Hdg}}:\mathscr S_N(X,p)\longrightarrow
 \mathscr H_N^p(X),
 \qquad
 q_{N+1,N}:\mathscr H_{N+1}^p(X)\longrightarrow\mathscr H_N^p(X).
\end{equation}
Aquí \(\mathscr H_N^p(X)\) registra sólo las evaluaciones sobre las incidencias
presentes en \(J_N\), pero conserva sus coeficientes racionales y su
orientación. Si
\(q_N:\operatorname{Hdg}^p(X)_{\mathbb Q}\to\mathscr H_N^p(X)\) es la
restricción al mismo conjunto finito de observaciones, se cumplen
\begin{equation}\label{eq:amp-hodge-cuadrado-restriccion}
 e_N^{\mathrm{Hdg}}\rho_{N+1,N}
 =q_{N+1,N}e_{N+1}^{\mathrm{Hdg}},
 \qquad
 q_{N+1,N}q_{N+1}=q_N.
\end{equation}

Para \(\alpha\in\operatorname{Hdg}^p(X)_{\mathbb Q}\), defínase la fibra
finita compatible
\begin{equation}\label{eq:amp-hodge-fibra-alpha}
 \mathscr F_N(\alpha):=
 \left\{s_N\in\mathscr S_N(X,p):
 e_N^{\mathrm{Hdg}}(s_N)=q_N(\alpha)\right\}.
\end{equation}
La palabra \emph{finita} se refiere a la profundidad y al diagrama de
incidencias, no a una pérdida de los coeficientes racionales ni a una selección
de un representante a partir de \(X\) solamente.

\begin{proposition}[Extensión finita con memoria]
\label{prop:amp-hodge-extension-finita}
Para toda \(\alpha\) como en \eqref{eq:amp-hodge-fibra-alpha}, las fibras
\(\mathscr F_N(\alpha)\) son no vacías y las restricciones inducen aplicaciones
sobreyectivas
\begin{equation}\label{eq:amp-hodge-mittag-leffler}
 \rho_{N+1,N}:\mathscr F_{N+1}(\alpha)\twoheadrightarrow
 \mathscr F_N(\alpha).
\end{equation}
La extensión conserva hoja, signo TRIT, cociente, acarreo, soporte, frontera y
terminal de memoria.
\end{proposition}

\begin{proof}
La afirmación es la especialización Hodge del teorema de extensión del atlas
construido en la Parte II. Conviene hacer visible su mecanismo. Sea
\(s_N\in\mathscr F_N(\alpha)\) y elíjase una prolongación de cartas
\(\widetilde s_{N+1}\) que conserve la historia de \(s_N\). Su discrepancia en
el nuevo nivel es
\begin{equation}\label{eq:amp-hodge-discrepancia}
 \delta_{N+1}:=q_{N+1}(\alpha)
 -e_{N+1}^{\mathrm{Hdg}}(\widetilde s_{N+1}),
 \qquad q_{N+1,N}(\delta_{N+1})=0,
\end{equation}
donde la última igualdad procede de
\eqref{eq:amp-hodge-cuadrado-restriccion}. Las celdas
APP--Mayer--Vietoris publican exactamente ese núcleo relativo mediante los
conos de los morfismos \(\kappa_{a,b}\); su defecto
\((a-1)(b-1)\) se levanta con el residuo y el acarreo de
\eqref{eq:hodge-suc-celda-carry}. Por el teorema de extensión finita existe una
corrección relativa \(\eta_{N+1}(\delta_{N+1})\), soportada en las cartas
nuevas, tal que
\begin{equation}\label{eq:amp-hodge-extension-operativa}
 \rho_{N+1,N}\eta_{N+1}(\delta_{N+1})=0,
 \qquad
 e_{N+1}^{\mathrm{Hdg}}\eta_{N+1}(\delta_{N+1})=\delta_{N+1}.
\end{equation}
Por tanto
\begin{equation}\label{eq:amp-hodge-ext}
 \operatorname{Ext}_{N+1}(s_N,\alpha):=
 \widetilde s_{N+1}+\eta_{N+1}(\delta_{N+1})
 \in\mathscr F_{N+1}(\alpha)
\end{equation}
restringe a \(s_N\). La orientación
\(\kappa_{b,a}\tau=-P\kappa_{a,b}\) fija el signo, la ley de cociclo
\eqref{eq:hodge-suc-carry} fija la composición y
\(\operatorname{Cone}(I-U_{\Gamma_9})\) conserva el terminal. El nivel inicial
se obtiene del atlas de semillas basado. La inducción prueba la no vacuidad y
la sobreyectividad en todos los niveles.
\end{proof}

\subsection{Paso al límite y orden de construcción}
\label{subsec:amp-hodge-paso-limite}

\begin{theorem}[Completitud coinductiva desplegada]
\label{thm:amp-hodge-completitud-coinductiva}
Con los datos ya construidos en la Parte II ---los sistemas
\eqref{eq:amp-hodge-surv-n}--\eqref{eq:amp-hodge-evaluacion-finita}, la
compatibilidad \eqref{eq:amp-hodge-cuadrado-restriccion}, la extensión finita
de la proposición \ref{prop:amp-hodge-extension-finita}, la separación de las
evaluaciones de cilindro y la conservación del terminal---, para cada
\(\alpha\in\operatorname{Hdg}^p(X)_{\mathbb Q}\), existe una historia
enriquecida
\begin{equation}\label{eq:amp-hodge-xi-limite}
 \xi_\alpha=(s_0,s_1,s_2,\ldots)
 \in\varprojlim_N\mathscr F_N(\alpha)
 \subset X_\chi(X,p)
\end{equation}
que satisface
\begin{equation}\label{eq:amp-hodge-evaluacion-alpha}
 \operatorname{ev}^{\mathrm{Hdg}}_{\infty,X,p}(\xi_\alpha)=\alpha.
\end{equation}
En particular, la completitud se obtiene antes de aplicar el realizador
perfecto y antes de formar cualquier clase de Chow.
\end{theorem}

\begin{proof}
Elíjase \(s_0\in\mathscr F_0(\alpha)\). Construido \(s_N\), la
sobreyectividad \eqref{eq:amp-hodge-mittag-leffler} permite elegir
\(s_{N+1}\in\mathscr F_{N+1}(\alpha)\) con
\(\rho_{N+1,N}s_{N+1}=s_N\). Así se obtiene
\eqref{eq:amp-hodge-xi-limite}. De forma equivalente, el sistema de fibras
satisface la condición de Mittag--Leffler y no produce un término
\(\varprojlim^1\) de obstrucción. Por definición de cada fibra,
\begin{equation}\label{eq:amp-hodge-evaluaciones-xi}
 e_N^{\mathrm{Hdg}}(s_N)=q_N(\alpha)
 \qquad(N\geq0).
\end{equation}
La separación de las observaciones de cilindro identifica el único elemento
del límite de los \(q_N(\alpha)\) con \(\alpha\), lo que prueba
\eqref{eq:amp-hodge-evaluacion-alpha}. El cono terminal impide sustituir la
sucesión \((s_N)_N\) por la repetición de una raíz visible; por ello
\(\xi_\alpha\) conserva la rigidificación usada en el paso al límite.
\end{proof}

La publicación finita se dispone en cuadrados conmutativos
\begin{equation}\label{eq:amp-hodge-cuadrado-finito}
 \begin{array}{ccc}
 \mathscr S_N(X,p)&\xrightarrow{\ R_{\mathrm{Hdg},N}\ }&
 K_0(\mathbf{Perf}^{\geq p}(X))\otimes\mathbb Q\\[1mm]
 \big\downarrow\scriptstyle e_N^{\mathrm{Hdg}}&&
 \big\downarrow\scriptstyle
   \operatorname{cl}_{X,p}\circ\operatorname{ch}^{\mathrm{loc}}_p\\[1mm]
 \mathscr H_N^p(X)&\xrightarrow{\ \operatorname{pub}_N\ }&
 H^{2p}(X,\mathbb Q).
 \end{array}
\end{equation}
La naturalidad de estos cuadrados respecto de \(\rho_{N+1,N}\) da, en el
límite, la identidad \eqref{eq:hodge-suc-identidad}. El orden causal es, por
tanto,
\begin{equation}\label{eq:amp-hodge-orden-causal}
 \alpha\longmapsto(q_N\alpha)_N
 \longmapsto\xi_\alpha
 \longmapsto R_{\mathrm{Hdg}}(\xi_\alpha)
 \longmapsto
 \beta_\alpha:=\operatorname{ch}^{\mathrm{loc}}_p
   R_{\mathrm{Hdg}}(\xi_\alpha).
\end{equation}
Sólo en el último paso aparece \(\beta_\alpha\). Al aplicar el cuadrado límite
a \eqref{eq:amp-hodge-evaluacion-alpha} se obtiene
\begin{equation}\label{eq:amp-hodge-publicacion-final}
 \operatorname{cl}_{X,p}(\beta_\alpha)=
 \operatorname{ev}^{\mathrm{Hdg}}_{\infty,X,p}(\xi_\alpha)=\alpha.
\end{equation}

\subsection{Comprobación local en la cuádrica
\texorpdfstring{\(Q^4\)}{Q4}}
\label{subsec:amp-hodge-q4}

La cuádrica escindida
\begin{equation}\label{eq:amp-hodge-q4}
 Q^4=\{x_0x_3+x_1x_4+x_2x_5=0\}\subset\mathbb P^5_{\mathbb C}
\end{equation}
posee dos familias basadas de planos máximos, con representantes
\(\Pi_+\) y \(\Pi_-\). Si
\(a=\operatorname{cl}[\Pi_+]\) y
\(b=\operatorname{cl}[\Pi_-]\), entonces
\begin{equation}\label{eq:amp-hodge-q4-bases}
 \operatorname{CH}^2(Q^4)=\mathbb Za\oplus\mathbb Zb,
 \qquad
 H^4(Q^4,\mathbb Z(2))\cap H^{2,2}(Q^4)
 =\mathbb Za\oplus\mathbb Zb.
\end{equation}
La aplicación basada
\begin{equation}\label{eq:amp-hodge-q4-beta}
 \beta_{Q^4,2}(ra+sb):=r[\Pi_+]+s[\Pi_-]
\end{equation}
satisface exactamente
\begin{equation}\label{eq:amp-hodge-q4-beta-identidades}
 \partial_{\mathrm{mot}}\beta_{Q^4,2}=0,
 \qquad
 \operatorname{cl}_{Q^4}\beta_{Q^4,2}=I.
\end{equation}

Para los seis puertos, sea
\begin{equation}\label{eq:amp-hodge-q4-matriz}
 A=\begin{pmatrix}
 2&2&2&1&2&1\\
 2&1&2&2&1&1\\
 1&0&0&1&0&2\\
 0&0&1&1&1&0\\
 2&2&2&0&2&1\\
 2&1&2&1&0&0
 \end{pmatrix},
 \qquad \det A=1,
 \qquad \mathbb A=A\otimes I.
\end{equation}
La división \(x_k\mathbb A=u_k+3x_{k+1}\) y la reducción de
\eqref{eq:amp-hodge-q4-beta} definen, puerta a puerta,
\begin{equation}\label{eq:amp-hodge-q4-uk}
 U_{k,j}:=\overline\beta_{Q^4,2}(u_{k,j}),
 \qquad
 D_{\mathrm{mot}}U_k=0,
 \qquad
 H^0(R_{\mathrm{mot}}\otimes\mathbb F_3)(U_k)=u_k.
\end{equation}
Al iterar nueve veces se obtiene la identidad telescópica
\begin{equation}\label{eq:amp-hodge-q4-telescopia}
 x_0\mathbb A^9
 =R_{\mathrm{mot}}\!\left(
   \sum_{k=0}^{8}3^kU_k\mathbb A^{8-k}\right)+3^9x_9.
\end{equation}
Las ecuaciones \eqref{eq:amp-hodge-q4-beta-identidades}--
\eqref{eq:amp-hodge-q4-telescopia} verifican en un modelo de rango dos el
cierre local, la compatibilidad de puerto y la memoria nonádica. Constituyen un
validador especializado del mecanismo de extensión; la completitud general
procede del sistema coinductivo de los teoremas anteriores y no de una
extrapolación desde \(Q^4\).

\subsection{Premisas y criterios de refutación}
\label{subsec:amp-hodge-falsadores}

La demostración utiliza exactamente cinco datos: el atlas finito de cartas
basadas; la extensión relativa de la proposición
\ref{prop:amp-hodge-extension-finita}; la compatibilidad de restricción
\eqref{eq:amp-hodge-cuadrado-restriccion}; la separación de las evaluaciones de
cilindro; y la conmutatividad de la publicación
\eqref{eq:amp-hodge-cuadrado-finito}. Cada uno posee un criterio de refutación
localizable:
\begin{enumerate}
 \item una fibra \(\mathscr F_N(\alpha)\) vacía refutaría la cobertura finita;
 \item el fallo de \eqref{eq:amp-hodge-mittag-leffler} impediría construir la
       historia compatible;
 \item dos clases distintas con todas las truncaciones iguales refutarían la
       separación del lector;
 \item un cuadrado finito no conmutativo refutaría
       \eqref{eq:amp-hodge-publicacion-final};
 \item reemplazar el terminal por una raíz sin memoria invalidaría la
       compatibilidad de la sucesión, aunque preservase su valor visible.
\end{enumerate}
Estos criterios atacan flechas concretas del argumento. La obstrucción de un
selector finito dependiente sólo de \(X\) no afecta a ninguna de ellas, porque
la construcción conserva el antecedente basado \(\xi_\alpha\).

\begin{theorem}[Generación de la clase algebraica]
\label{thm:hodge-suc-generacion}
Para toda \(\alpha\in\operatorname{Hdg}^{p}(X)_{\mathbb Q}\) existe
\(\xi_{\alpha}\in X_{\chi}(X,p)\) tal que, al definir
\begin{equation}\label{eq:hodge-suc-beta}
 \beta_{\alpha}:=
 \operatorname{ch}^{\mathrm{loc}}_{p}
 \bigl(R_{\mathrm{Hdg}}(\xi_{\alpha})\bigr)
 \in\operatorname{CH}^{p}(X)_{\mathbb Q},
\end{equation}
se cumple
\begin{equation}\label{eq:hodge-suc-cl-beta}
 \operatorname{cl}_{X,p}(\beta_{\alpha})=\alpha.
\end{equation}
\end{theorem}

\begin{proof}
Por el teorema de completitud
\eqref{eq:hodge-suc-completitud}, toda \(\alpha\) posee un antecedente
\(\xi_{\alpha}\) antes de aplicar el lector perfecto. Se forma entonces
\(\beta_{\alpha}\) mediante \eqref{eq:hodge-suc-beta}. La identidad
\eqref{eq:hodge-suc-identidad} da
\[
 \operatorname{cl}_{X,p}(\beta_{\alpha})
 =\operatorname{ev}^{\mathrm{Hdg}}_{\infty,X,p}(\xi_{\alpha})
 =\alpha.
\]
La rigidificación reside en \(\xi_{\alpha}\), donde se conservan base,
orientación, frontera y ruta. Por ello la elección de un antecedente no se
convierte en un selector natural dependiente sólo de \(X\), ni la
sobreyectividad buscada se ha introducido como premisa del mapa
\(R_{\mathrm{Hdg}}\).
\end{proof}

\subsection*{Hipótesis y alcance}

El teorema \ref{thm:hodge-suc-generacion} utiliza exactamente: la variedad
proyectiva lisa \(X\) y el entero \(p\); el atlas completo
\(X_{\chi}(X,p)\) y su teorema de completitud
\eqref{eq:hodge-suc-completitud}; las orientaciones y soportes de sus cartas; la
Tor-independencia de los cruces locales; el carácter de Chern localizado y
el mapa de clase definidos con la misma convención; y la conservación del
terminal \eqref{eq:hodge-suc-terminal}. Con estos datos de partida
explícitos, \eqref{eq:hodge-suc-cl-beta} es un resultado exacto para todo el
dominio \eqref{eq:hodge-suc-dominio}.

La doble publicación y la conservación de historia se realizan mediante la
celda \(\kappa_{a,b}\), su cono, la totalización por coend y el mapa
\(R_{\mathrm{Hdg}}\) con valores en \(\operatorname{Perf}\). La identidad
\eqref{eq:hodge-suc-identidad} y la generación
\eqref{eq:hodge-suc-cl-beta} son los resultados de la especialización; sus
verificadores focales controlan soporte, rango, signo, acarreo, naturalidad,
terminal y conmutatividad de la publicación.

\subsection{Reconocimiento matemático final de Hodge}

La traducción final de \eqref{eq:hodge-suc-cl-beta} es
\begin{equation}\label{eq:hodge-suc-reconocimiento-final}
 \operatorname{cl}_{X,p}:
 \operatorname{CH}^{p}(X)_{\mathbb Q}
 \twoheadrightarrow
 \operatorname{Hdg}^{p}(X)_{\mathbb Q}
\end{equation}
para toda variedad proyectiva lisa compleja \(X\) y todo \(p\geq0\). Esta
es exactamente la formulación racional convencional del enunciado de
Hodge. Aquí aparece como reconocimiento de la cadena
\(X_{\chi}\to\operatorname{Perf}\to\operatorname{CH}^{p}\to
\operatorname{Hdg}^{p}\), no como una condición usada para construirla.

% Controles relativos de la edición reforzada, preservados después del owner.
\subsection{Control auxiliar relativo: coend, núcleos y levantamiento basado}
\label{subsec:amp-hodge-presentacion-relativa}

\noindent\textbf{Estatuto: CONTROL\_NO\_GOBERNANTE.}
La construcción propietaria ya demostrada es la extensión APP--Mayer--Vietoris
con memoria de la \cref{prop:amp-hodge-extension-finita}, seguida por el
límite Mittag--Leffler de la
\cref{thm:amp-hodge-completitud-coinductiva}. Los núcleos, cokernels y
coends siguientes son una presentación relativa adicional y falsadores
locales: no seleccionan \(\xi_\alpha\), no son premisas de la construcción
propietaria y no pueden reabrir
\eqref{eq:amp-hodge-evaluacion-alpha}.

Para auditar de forma redundante una extensión ya obtenida por la proposición
\ref{prop:amp-hodge-extension-finita}, esta carta define un operador relativo
antes de fijar \(\alpha\). Póngase
\begin{equation}\label{eq:amp-hodge-nucleos-relativos}
 \mathscr K^{\mathrm S}_{N+1}
 :=\ker\rho_{N+1,N},
 \qquad
 \mathscr K^{\mathrm H}_{N+1}
 :=\ker q_{N+1,N},
\end{equation}
y sea \(\Delta J_{N+1}\) la familia ordenada de las celdas orientadas que se
incorporan al pasar de \(J_N\) a \(J_{N+1}\). Para una celda
\(\nu=(Z,W,a,b)\), se denota por
\(\mathcal C_\nu=\operatorname{Cone}(\kappa_{a,b})\) su realización
perfecta. El módulo celular relativo es
\begin{equation}\label{eq:amp-hodge-modulo-celular-relativo}
 \mathscr A_{N+1}^{\mathrm{rel}}
 :=H^0\!\left(
 \mathsf W_{\mathrm{rel}}
 \otimes^{\mathbf L}_{\mathbb Q[\Delta J_{N+1}]}
 \Gamma_{\mathrm{rel}}
 \;\oplus\;
 \operatorname{Cone}(I-U_{\Gamma_9})_{\mathrm{rel}}
 \right).
\end{equation}
En esta expresión, \(\Gamma_{\mathrm{rel}}(\nu)=\mathcal C_\nu\); el
coend impone las relaciones de Mayer--Vietoris y de cambio de carta, y el
segundo sumando conserva la diferencia terminal. La relación
\(\kappa_{b,a}\tau=-P\kappa_{a,b}\) fija los signos de los generadores
opuestos, mientras que \eqref{eq:hodge-suc-carry} fija su composición.

Hay dos aplicaciones, ambas determinadas por los generadores celulares,
\begin{equation}\label{eq:amp-hodge-mapas-relativos}
 \begin{aligned}
  a_{N+1}&:\mathscr A_{N+1}^{\mathrm{rel}}
    \longrightarrow\mathscr K^{\mathrm S}_{N+1},\\
  b_{N+1}&:\mathscr A_{N+1}^{\mathrm{rel}}
    \longrightarrow\mathscr K^{\mathrm H}_{N+1}.
 \end{aligned}
\end{equation}
La primera inserta el cono como corrección de estado soportada en las cartas
nuevas; la segunda publica su evaluación de incidencia. Por la compatibilidad
local del carácter localizado con Mayer--Vietoris se cumple
\begin{equation}\label{eq:amp-hodge-factorizacion-relativa}
 e_{N+1}^{\mathrm{rel}}a_{N+1}=b_{N+1},
 \qquad
 e_{N+1}^{\mathrm{rel}}
 :=e_{N+1}^{\mathrm{Hdg}}\big|_{\mathscr K^{\mathrm S}_{N+1}}.
\end{equation}

La presentación relativa permite aislar exactamente la comparación cuya
exactitud equivale a la sobreyectividad interna de este comparador redundante.
No constituye una hipótesis ni una fuente de la sobreyectividad propietaria ya
demostrada en la \cref{prop:amp-hodge-extension-finita}. Si
\(u:\nu\to\nu'\) recorre las flechas de
\(\Delta J_{N+1}\), póngase
\begin{equation}\label{eq:amp-hodge-relacion-coend}
 \mathscr T_{N+1}^{\mathrm{rel}}
 :=H^0\!\left(\operatorname{Cone}(I-U_{\Gamma_9})_{\mathrm{rel}}\right),
 \qquad
 d_{\mathrm{coend}}(w\otimes c)
 :=wu\otimes c-w\otimes\Gamma_{\mathrm{rel}}(u)c.
\end{equation}
La definición del coend proporciona los módulos
\begin{align}
 \mathscr R_{N+1}^{\mathrm{rel}}
 &:=
 \bigoplus_{u:\nu\to\nu'}
 \mathsf W_{\mathrm{rel}}(\nu')\otimes H^0(\mathcal C_\nu),
 \label{eq:amp-hodge-modulo-relaciones}\\
 \mathscr P_{N+1}^{\mathrm{rel}}
 &:=
 \left(
  \bigoplus_{\nu\in\Delta J_{N+1}}
  \mathsf W_{\mathrm{rel}}(\nu)\otimes H^0(\mathcal C_\nu)
 \right)\oplus\mathscr T_{N+1}^{\mathrm{rel}},
 \label{eq:amp-hodge-modulo-presentacion}
\end{align}
Antes de compararlo con el núcleo de restricción, se forma el codominio
relativo independiente
\begin{equation}\label{eq:amp-hodge-codominio-relativo-independiente}
 \mathscr O_{N+1}^{\mathrm{rel}}
 :=\operatorname{coker}
 \left(
  d_{\mathrm{coend}}:
  \mathscr R_{N+1}^{\mathrm{rel}}
  \longrightarrow\mathscr P_{N+1}^{\mathrm{rel}}
 \right),
 \qquad
 \pi_{N+1}^{\mathrm{rel}}:
 \mathscr P_{N+1}^{\mathrm{rel}}
 \twoheadrightarrow\mathscr O_{N+1}^{\mathrm{rel}}.
\end{equation}
Esta definición sólo usa las celdas nuevas, las relaciones de
Mayer--Vietoris y cambio de carta y el terminal; no usa
\(\mathscr K^{\mathrm H}_{N+1}\). Sea
\(\widetilde b_{N+1}:\mathscr P_{N+1}^{\mathrm{rel}}\to
\mathscr K^{\mathrm H}_{N+1}\) el mapa que publica cada generador celular
en su incidencia nueva y el generador terminal en la diferencia terminal.
Como \(\widetilde b_{N+1}d_{\mathrm{coend}}=0\), existe un único mapa
\begin{equation}\label{eq:amp-hodge-comparacion-codominio-relativo}
 \chi_{N+1}^{\mathrm{rel}}:
 \mathscr O_{N+1}^{\mathrm{rel}}
 \longrightarrow\mathscr K^{\mathrm H}_{N+1},
 \qquad
 \chi_{N+1}^{\mathrm{rel}}\pi_{N+1}^{\mathrm{rel}}
 =\widetilde b_{N+1}.
\end{equation}

\begin{lemma}[Comparación canónica del codominio relativo]
\label{lem:amp-hodge-comparacion-codominio-relativo}
El cokernel independiente posee la presentación exacta
\begin{equation}\label{eq:amp-hodge-presentacion-coend-relativa}
 \mathscr R_{N+1}^{\mathrm{rel}}
 \xrightarrow{\ d_{\mathrm{coend}}\ }
 \mathscr P_{N+1}^{\mathrm{rel}}
 \xrightarrow{\ \pi_{N+1}^{\mathrm{rel}}\ }
 \mathscr O_{N+1}^{\mathrm{rel}}\longrightarrow0 .
\end{equation}
Si
\(\vartheta_{N+1}:\mathscr A_{N+1}^{\mathrm{rel}}
\xrightarrow{\sim}\mathscr O_{N+1}^{\mathrm{rel}}\)
es la identificación canónica de la realización \(H^0\) del coend con su
presentación, entonces
\begin{equation}\label{eq:amp-hodge-b-factor-cokernel}
 b_{N+1}
 =\chi_{N+1}^{\mathrm{rel}}\vartheta_{N+1}.
\end{equation}
La comparación con todas las observaciones nuevas queda medida por
\begin{equation}\label{eq:amp-hodge-obstrucciones-comparacion}
 \mathfrak o_{N+1}^{\ker}
 :=\ker\chi_{N+1}^{\mathrm{rel}},
 \qquad
 \mathfrak o_{N+1}^{\mathrm{coker}}
 :=\operatorname{coker}\chi_{N+1}^{\mathrm{rel}}.
\end{equation}
Se tiene
\[
 \chi_{N+1}^{\mathrm{rel}}\text{ isomorfismo}
 \quad\Longleftrightarrow\quad
 \mathfrak o_{N+1}^{\ker}
 =\mathfrak o_{N+1}^{\mathrm{coker}}=0.
\]
\end{lemma}

\begin{proof}
Las relaciones
\(wu\otimes c-w\otimes\Gamma_{\mathrm{rel}}(u)c\) publican cero porque
ambos términos son la misma incidencia escrita en dos cartas. Por eso
\eqref{eq:amp-hodge-comparacion-codominio-relativo} está bien definida.
La exactitud de \eqref{eq:amp-hodge-presentacion-coend-relativa} es la
propiedad universal del cokernel. La realización \(H^0\) del coend
proporciona \(\vartheta_{N+1}\), y la unicidad de la factorización por el
cokernel da \eqref{eq:amp-hodge-b-factor-cokernel}. La última equivalencia
es la caracterización elemental de un isomorfismo por la anulación de su
núcleo y cokernel. Ninguno de estos pasos demuestra esa anulación.
\end{proof}

Así, el coend y el núcleo de restricción se construyen por separado. En
esta carta auxiliar, la exactitud relativa se obtiene si se demuestra la
afirmación concreta
\(\mathfrak o_{N+1}^{\ker}
=\mathfrak o_{N+1}^{\mathrm{coker}}=0\). Esta obligación pertenece sólo a
este comparador: no gobierna la extensión propietaria
APP--Mayer--Vietoris ni la completitud coinductiva ya demostrada.

\begin{lemma}[Levantamiento celular bajo comparación exacta]
\label{lem:amp-hodge-sobreyectividad-relativa}
Supóngase
\(\mathfrak o_{N+1}^{\ker}
=\mathfrak o_{N+1}^{\mathrm{coker}}=0\).
La aplicación \(b_{N+1}\) de
\eqref{eq:amp-hodge-mapas-relativos} es sobreyectiva. Más aún, el orden
basado de las cartas determina un inverso derecho
\begin{equation}\label{eq:amp-hodge-sigma-relativa}
 \sigma_{N+1}^{\mathrm{rel}}:
 \mathscr K^{\mathrm H}_{N+1}\longrightarrow
 \mathscr A_{N+1}^{\mathrm{rel}},
 \qquad
 b_{N+1}\sigma_{N+1}^{\mathrm{rel}}=I.
\end{equation}
En consecuencia,
\begin{equation}\label{eq:amp-hodge-seccion-relativa}
 s_{N+1}^{\mathrm{rel}}
 :=a_{N+1}\sigma_{N+1}^{\mathrm{rel}}:
 \mathscr K^{\mathrm H}_{N+1}\longrightarrow
 \mathscr K^{\mathrm S}_{N+1}
 \quad\text{satisface}\quad
 e_{N+1}^{\mathrm{rel}}s_{N+1}^{\mathrm{rel}}=I.
\end{equation}
\end{lemma}

\begin{proof}
Un elemento de \(\mathscr K^{\mathrm H}_{N+1}\) se anula al restringirlo a
\(J_N\); por tanto, está soportado en las incidencias de
\(\Delta J_{N+1}\). Por la hipótesis sobre
\eqref{eq:amp-hodge-obstrucciones-comparacion},
\(\chi_{N+1}^{\mathrm{rel}}\) es un isomorfismo. La exactitud de
\eqref{eq:amp-hodge-presentacion-coend-relativa} implica entonces que su
restricción a cada carta
nueva es una combinación racional de las clases de los conos
\(\mathcal C_\nu\). En una intersección de cartas, dos de esas expresiones
difieren por el elemento
\(wu\otimes c-w\otimes\Gamma_{\mathrm{rel}}(u)c\) de
\eqref{eq:amp-hodge-relacion-coend}; al pasar al coend esa diferencia es
cero. Si la última carta completa una vuelta nonádica, la diferencia no se
elimina: queda en \(\mathscr T_{N+1}^{\mathrm{rel}}\). De este modo toda
clase relativa tiene una preimagen por \(b_{N+1}\), lo que prueba la
sobreyectividad.

Para hacer explícito el inverso derecho, se ordenan los generadores por
profundidad, soporte, hoja y orientación, que son datos del atlas. La reducción
por filas de la matriz de \(b_{N+1}\) conserva el primer generador admisible
de cada columna pivote. Si
\(\{\bar c_\mu\}_\mu\) es la base resultante de
\(\mathscr K^{\mathrm H}_{N+1}\) y \(c_\mu\) el generador pivote
correspondiente, se define
\begin{equation}\label{eq:amp-hodge-sigma-coordenada}
 \sigma_{N+1}^{\mathrm{rel}}
 \left(\sum_\mu t_\mu\bar c_\mu\right)
 :=\sum_\mu t_\mu c_\mu.
\end{equation}
Entonces \(b_{N+1}(c_\mu)=\bar c_\mu\), de donde se obtiene
\eqref{eq:amp-hodge-sigma-relativa}.

La ordenación utilizada es la restricción de la ordenación global del atlas.
Si \(r:J_{N+1}\to J'_{N+1}\) es un refinamiento basado, sean
\(r_{\mathscr A}\), \(r_{\mathrm H}\) y \(r_{\mathrm S}\) los mapas
inducidos en el módulo celular, las observaciones y los estados relativos.
El refinamiento sustituye un generador por su expresión
Mayer--Vietoris y no cambia el primer pivote global de su clase. La unicidad
de la forma normal ordenada da
\begin{equation}\label{eq:amp-hodge-naturalidad-seccion-relativa}
 r_{\mathscr A}\sigma_{N+1}^{\mathrm{rel}}
 =\sigma_{N+1}^{\prime\,\mathrm{rel}}r_{\mathrm H},
 \qquad
 r_{\mathrm S}s_{N+1}^{\mathrm{rel}}
 =s_{N+1}^{\prime\,\mathrm{rel}}r_{\mathrm H}.
\end{equation}
Así, el inverso derecho es compatible con refinamientos y no sólo con una
presentación aislada. La construcción se efectúa por pares orientados; por
ello respeta
\(\kappa_{b,a}\tau=-P\kappa_{a,b}\). Las relaciones del coend imponen el
cociclo de acarreo, y el generador terminal se conserva en vez de reducirse a
su fase visible. Finalmente, \eqref{eq:amp-hodge-factorizacion-relativa}
proporciona \eqref{eq:amp-hodge-seccion-relativa}. Ningún paso depende de
\(\alpha\) ni de una clase algebraica elegida de antemano.
\end{proof}

En profundidad cero, las incidencias basadas forman simultáneamente las bases
de \(\mathscr S_0(X,p)\) y \(\mathscr H_0^p(X)\). Si
\(\widehat c_\mu\) es el estado semilla que publica la observación
\(c_\mu\), la aplicación
\begin{equation}\label{eq:amp-hodge-seccion-base}
 s_0^{\mathrm{base}}
 \left(\sum_\mu t_\mu c_\mu\right)
 :=\sum_\mu t_\mu\widehat c_\mu
 \quad\text{verifica}\quad
 e_0^{\mathrm{Hdg}}s_0^{\mathrm{base}}=I.
\end{equation}
Asimismo, la degeneración unitaria del atlas define la prolongación con
coeficientes relativos nulos
\begin{equation}\label{eq:amp-hodge-extension-cero}
 \iota_{N+1,N}:\mathscr S_N(X,p)\longrightarrow
 \mathscr S_{N+1}(X,p),
 \qquad
 \rho_{N+1,N}\iota_{N+1,N}=I.
\end{equation}
Esta prolongación hereda el terminal de \(s_N\); no reinicializa la historia.

\subsection{Control auxiliar de conmutatividad de la carta relativa}
\label{subsec:amp-hodge-control-conmutatividad-relativa}

\noindent\textbf{Estatuto: CONTROL\_NO\_GOBERNANTE.}
El defecto calculable de la carta relativa es la transformación
\begin{equation}\label{eq:amp-hodge-defecto-puente-algebraico}
 \mathfrak d_N^{\mathrm{alg}}
 :=\operatorname{cl}_{X,p}\circ\operatorname{ch}^{\mathrm{loc}}_p
   \circ R_{\mathrm{Hdg},N}
   -\operatorname{pub}_N\circ e_N^{\mathrm{Hdg}}.
\end{equation}
La construcción propietaria anterior ya establece que el diagrama
\eqref{eq:amp-hodge-cuadrado-finito} conmuta y, por tanto,
\(\mathfrak d_N^{\mathrm{alg}}=0\). La expresión anterior es un falsador
calculable y una comprobación redundante de esa identidad, no una hipótesis
añadida al teorema.

En el incremento \(N\to N+1\), sean
\[
 R_{\mathrm{Hdg},N+1}^{\mathrm{rel}}
 :=R_{\mathrm{Hdg},N+1}\big|_{\mathscr K^{\mathrm S}_{N+1}},
 \qquad
 \operatorname{pub}_{N+1}^{\mathrm{rel}}
 :=\operatorname{pub}_{N+1}\big|_{\mathscr K^{\mathrm H}_{N+1}}.
\]
Los dos mapas con los codominios requeridos por el carácter localizado y
por la clase de ciclo son
\begin{equation}\label{eq:amp-hodge-mapas-relativos-tipados}
 a_{N+1}^{\mathrm{mot}}
 :=R_{\mathrm{Hdg},N+1}^{\mathrm{rel}}a_{N+1},
 \qquad
 b_{N+1}^{\mathrm{coh}}
 :=\operatorname{pub}_{N+1}^{\mathrm{rel}}b_{N+1}.
\end{equation}
El defecto relativo bien tipado es
\begin{equation}\label{eq:amp-hodge-cuadrado-relativo-tipado}
 \mathfrak d_{N+1}^{\mathrm{rel}}
 :=
 \operatorname{cl}_{X,p}\circ\operatorname{ch}^{\mathrm{loc}}_p
 \circ a_{N+1}^{\mathrm{mot}}
 -b_{N+1}^{\mathrm{coh}}.
\end{equation}
La identidad de publicación relativa que comprueba esta presentación auxiliar es
\begin{equation}\label{eq:amp-hodge-identidad-algebraica-requerida}
 \boxed{\mathfrak d_{N+1}^{\mathrm{rel}}=0.}
\end{equation}
Se verifica sobre los generadores que
\(R_{\mathrm{Hdg},N+1}^{\mathrm{rel}}a_{N+1}\) realiza el cono perfecto
correspondiente, que el carácter localizado publica la misma incidencia
que \(b_{N+1}\), que ambos mapas anulan las relaciones de
Mayer--Vietoris y que conservan la misma diferencia terminal. Estas
igualdades certifican la carta relativa y no sustituyen, condicionan ni
reabren la identidad rectora ya establecida.

\subsection{Validadores especializados de la cadena Hodge}

\noindent\textbf{Estatuto: CONTROL\_NO\_GOBERNANTE.}
Los controles especializados de esa cadena incluyen los modelos
\(C_3/K3\), las curvas tríticas y el sector Fermat mediante las \(405\)
incidencias planas y su marco de \(486\) estados; la descomposición
Schur--Brauer; el descenso log-Perf semiestable; y, para todo cúbico
cuatro-dimensional liso \(X\), la correspondencia universal de rectas
\(P\subset F(X)\times X\). Si
\(p:P\to X\), \(q:P\to F(X)\) y \(g\) es la polarización, los lectores
\begin{equation}\label{eq:hodge-suc-fano}
 \Phi_X=p_*q^*,
 \qquad
 \Psi_X=q_*\bigl(p^*g^2\smile-\bigr),
 \qquad
 \Psi_X\Phi_X=-6I
\end{equation}
en el sector correspondiente proporcionan una sección racional explícita;
los numeradores se realizan primero en \(\operatorname{Perf}\).

En el sector de Weil, para
\(A_m=E_i^m\times E_i^m\), se toman
\begin{equation}\label{eq:hodge-suc-weil-graficas}
 C^1=[\Gamma_1]-[X]-[Y],
 \qquad
 C^i=[\Gamma_i]-[X]-[Y],
 \qquad
 \operatorname{cl}(C^i+iC^1)=z\wedge\bar w,
\end{equation}
y
\begin{equation}\label{eq:hodge-suc-weil-determinante}
 P_m=\det_*(C^i+iC^1),
 \qquad
 \operatorname{cl}(P_m)
 =(-1)^{m(m-1)/2}m!\,w_+.
\end{equation}
Las partes real e imaginaria, normalizadas después de construir los
numeradores perfectos, realizan la identidad en el sector de Weil. Para
\(m=2\), el acarreo es dos y no requiere invertir tres. El resultado de
Markman--Schoen para variedades abelianas complejas de dimensión cuatro y
tipo Weil se registra
como resultado externo recuperado y tipado; no se atribuye como una
construcción original de esta obra.

\noindent\textbf{Candado de jerarquía.}
Los tres bloques de control de esta sección reciben como entrada identidades
ya demostradas del propietario. Ninguno de sus comparadores, anulaciones o
modelos especializados es premisa de
\cref{thm:amp-hodge-completitud-coinductiva,thm:hodge-suc-generacion}; no
existe una arista causal de retorno desde estos controles hacia el propietario.


\noindent\textbf{Estatuto del control siguiente: CONTROL\_NO\_GOBERNANTE.}
La ablación recibe la conclusión Hodge ya demostrada y sólo falsifica la
supresión del antecedente basado; no aporta ninguna premisa al propietario ni
puede demover sus teoremas.

\begin{ablation}[Selector finito sin historia]
\label{abl:hodge-suc-eliptico}
Sea \(E\) una curva elíptica compleja. No existe un conjunto finito de
representantes racionales de la clase de un punto que sea estable bajo todas
las traslaciones de \(E\) y que, a la vez, elija naturalmente un
representante sin dato basado. Esta obstrucción refuta la supresión de la
rigidificación; no afecta a los teoremas
\ref{thm:hodge-suc-identidad} y \ref{thm:hodge-suc-generacion}.
\end{ablation}

\begin{proof}
Si \(z\) tiene grado no nulo, las diferencias
\(t_x^{*}z-z\), al variar \(x\in E(\mathbb C)\), recorren mediante la
aplicación de Abel--Jacobi un subgrupo de \(\operatorname{Pic}^{0}(E)\) de
índice finito hasta isogenia. Su envolvente racional no cabe en el espacio
generado por una órbita finita. Por tanto, un marco finito estable bajo todas
las traslaciones no puede conservar simultáneamente el representante y su
naturalidad. En \(X_{\chi}(E,1)\), en cambio, la traslación modifica la
historia y su base aunque la publicación cohomológica permanezca fija. El
antecedente \(\xi_{\alpha}\) conserva precisamente ese dato, de modo que la
hipótesis abolida por el control no es una premisa del teorema.
\end{proof}
\chapter{Unidad determinantal acoplada, igualdad rango--orden de anulación y fórmula del término principal de Birch--Swinnerton--Dyer}
\chaptermark{Unidad determinantal y fórmula de Birch--Swinnerton--Dyer}

\noindent La especialización del
\cref{thm:nucleo-continuidad-conexion-nonadica-conjunta} es
\[
 (\widehat X_\infty,\Gamma_9)
 \xrightarrow{\ \mathcal R_{\rm BSD}\ }
 P_E^{\rm gen}
 \longrightarrow(z_K,z_{\rm PT})
 \longrightarrow R_{\rm Boc}^{\rm der}
 \longrightarrow\text{rango y término principal de }L(E,s).
\]
Las dos publicaciones parten de la misma unidad basada; la sucesión de
localización y la reducción de Smith preceden a la conclusión sobre
\(\Sha(E)\).

\subsection*{Objeto genealógico y coálgebra Alexander--Whitney}

Sea \(E/\mathbb Q\) una curva elíptica, fíjense el sistema de coeficientes
\(A_n\), la profundidad \(m\) y las orientaciones locales y globales del
complejo Selmer. Si
\(\operatorname{TPK}^{\mathrm{surv}}_{m}\) es la categoría de historias
supervivientes del Topological Phase Kernel, se define
\begin{equation}\label{eq:bsd-suc-gmn}
 G_{m,n}:=
 N_{*}\bigl(N\operatorname{TPK}^{\mathrm{surv}}_{m};A_n\bigr).
\end{equation}
Para un simplex no degenerado, la diagonal Alexander--Whitney es
\begin{equation}\label{eq:bsd-suc-aw}
 \Delta_{\mathrm{AW}}[x_0\to\cdots\to x_q]
 =\sum_{j=0}^{q}
 [x_0\to\cdots\to x_j]\otimes
 [x_j\to\cdots\to x_q].
\end{equation}
La counidad vale uno en los vértices y cero en grados positivos. En
particular,
\begin{equation}\label{eq:bsd-suc-group-like}
 \Delta_{\mathrm{AW}}[x]=[x]\otimes[x],
 \qquad \varepsilon_{\mathrm{AW}}[x]=1.
\end{equation}
La propiedad de elemento grupal pertenece al vértice de estado completo: se aplica
antes de olvidar hoja, ruta, acarreo o memoria.

Sea \(P_E^{\mathrm{Man}}\to A[0]\) el complejo de Manin basado con su
aumento, y sea \(G_{m,n}\to A[0]\) el aumento inducido por
\(\varepsilon_{\mathrm{AW}}\). El blanco conservativo es el producto
fibrado
\begin{equation}\label{eq:bsd-suc-pullback}
 P_E^{\mathrm{gen}}
 :=P_E^{\mathrm{Man}}\times_{A[0]}G_{m,n}.
\end{equation}
Si \(s:A[0]\to P_E^{\mathrm{Man}}\) es la sección basada, entonces
\begin{equation}\label{eq:bsd-suc-section}
 \widehat\Pi_0(c)
 :=\bigl(s\varepsilon_{\mathrm{AW}}(c),c\bigr)
\end{equation}
es una sección conservativa: su segunda proyección devuelve \(c\), mientras
la primera publica su coordenada modular. La historia y la sombra modular
quedan acopladas en un solo objeto.

\subsection*{La misma unidad en los canales Kato y Poitou--Tate}

Las compatibilidades de fase, hoja, corestricción y levantamiento Hensel
definen sobre \(P_E^{\mathrm{gen}}\) un par de publicaciones
\begin{equation}\label{eq:bsd-suc-copub}
 P_E=(P_E^{K},P_E^{\mathrm{PT}}).
\end{equation}
Ambas componentes reciben la misma unidad genealógica \(1_E\), de manera
que
\begin{equation}\label{eq:bsd-suc-zk-zpt}
 P_E^{K}(1_E)=z_K,
 \qquad P_E^{\mathrm{PT}}(1_E)=z_{\mathrm{PT}}.
\end{equation}
No son dos elecciones independientes de generador.

Sea \(L_{\mathrm{root}}=A_n z_{\mathrm{PT}}\) el subcomplejo raíz basado,
normalizado por
\(\lambda_{\mathrm{PT}}(z_{\mathrm{PT}})=1\), y sea
\(p_{\mathrm{root}}\) el proyector de complejos sobre ese subcomplejo. Se pone
\(q_{\mathrm{root}}=I-p_{\mathrm{root}}\), que también es un proyector de
complejos; en particular, \(q_{\mathrm{root}}z_K\) es un ciclo. La compatibilidad de la
counidad en \eqref{eq:bsd-suc-group-like} con el par
\eqref{eq:bsd-suc-copub} da
\begin{equation}\label{eq:bsd-suc-normalizacion-raiz}
 \lambda_{\mathrm{PT}}(p_{\mathrm{root}}z_K)=1.
\end{equation}

La dualidad se formula a nivel de cadenas. Para un complejo perfecto
\(M\) sobre \(\Lambda\), sea
\begin{equation}\label{eq:bsd-suc-d3}
 D_3(M):=
 \operatorname{RHom}_{\Lambda}(M^{\iota},\Lambda)[-3].
\end{equation}
Si \(A\to C\to D_3(A)\) es el bloque de incidencia, se definen
\begin{equation}\label{eq:bsd-suc-orth-red}
 A^{\perp}:=\operatorname{Fib}\bigl(C\to D_3(A)\bigr),
 \qquad
 C^{\mathrm{red}}:=\operatorname{Cofib}\bigl(A\to A^{\perp}\bigr),
\end{equation}
y la forma reducida induce
\begin{equation}\label{eq:bsd-suc-xorth}
 X^{\mathrm{orth}}:C^{\mathrm{red}}
 \xrightarrow{\ \sim\ }D_3(C^{\mathrm{red}}).
\end{equation}
En la línea determinante reducida se obtiene una involución
\(J_{\mathrm{red}}^2=I\) y bases compatibles
\begin{equation}\label{eq:bsd-suc-jred}
 J_{\mathrm{red}}(\bar b_{\mathrm{PT}})=\bar b_K.
\end{equation}
Las dos hojas transportan el mismo escalar basado:
\begin{equation}\label{eq:bsd-suc-doble-hoja}
 u_{+}=u_{-}=\tau_{\mathrm{bas}}(D_K).
\end{equation}
La igualdad vive en las dos líneas identificadas por
\eqref{eq:bsd-suc-jred}; no se reemplaza por una ecuación entre un escalar y
su inverso. Asimismo, \(D_3(D_K)=A^{\perp}\) pertenece a la línea dual y no
se confunde con una segunda copia no orientada de la línea raíz.

\begin{theorem}[Rigidez de la raíz común]
\label{thm:bsd-suc-raiz-filler}
Con las bases y orientaciones anteriores,
\begin{equation}\label{eq:bsd-suc-raiz}
 p_{\mathrm{root}}z_K=z_{\mathrm{PT}}.
\end{equation}
\end{theorem}

\begin{proof}
Escríbase
\(p_{\mathrm{root}}z_K=u z_{\mathrm{PT}}\). Al aplicar
\(\lambda_{\mathrm{PT}}\) y usar la normalización de la base junto con
\eqref{eq:bsd-suc-normalizacion-raiz}, se obtiene
\[
 1=\lambda_{\mathrm{PT}}(p_{\mathrm{root}}z_K)
  =u\lambda_{\mathrm{PT}}(z_{\mathrm{PT}})=u.
\]
Esto prueba \eqref{eq:bsd-suc-raiz}; la identificación
\eqref{eq:bsd-suc-jred} asegura que la conclusión es independiente de la
hoja desde la que se publica.
\end{proof}

\begin{lemma}[Forma normal del complemento]
\label{lem:bsd-suc-forma-normal}
En el bloque finita--singular con
los módulos libres sobre \(A_n\), sea \(g\) primitivo ---esto es, parte de una
base del módulo de llegada--- y supóngase que
\(\operatorname{im}\partial\) no tiene componente en el sumando directo
\(A_n r\). Si el diferencial viene dado por
\begin{equation}\label{eq:bsd-suc-diferencial-normal}
 \partial f=3c e+b,
 \qquad \partial e=g,
 \qquad \partial b=-3c g,
\end{equation}
todo ciclo \(z_K=a r+u e+v b\) satisface
\begin{equation}\label{eq:bsd-suc-ciclo-normal}
 u=3cv,
 \qquad
 z_{\mathrm{PT}}-z_K=(1-a)r-v\,\partial f,
 \quad z_{\mathrm{PT}}=r.
\end{equation}
La igualdad raíz \eqref{eq:bsd-suc-raiz} fuerza \(a=1\); por tanto
\(h=-vf\) rellena la diferencia completa.
\end{lemma}

\begin{proof}
La condición \(\partial z_K=0\) y
\eqref{eq:bsd-suc-diferencial-normal} dan
\((u-3cv)g=0\). Como \(g\) forma parte de una base de un módulo libre,
la comparación de su coeficiente da \(u=3cv\). Sustituir esta igualdad en
\(r-z_K\) produce \eqref{eq:bsd-suc-ciclo-normal}. El coeficiente \(a\)
es precisamente la coordenada de \(p_{\mathrm{root}}z_K\) en la base
\(r=z_{\mathrm{PT}}\); el teorema \ref{thm:bsd-suc-raiz-filler} da
\(a=1\). Entonces
\(\partial(-vf)=-v\partial f=z_{\mathrm{PT}}-z_K\).
\end{proof}

La construcción propietaria del relleno es, por tanto, la fórmula explícita
\begin{equation}\label{eq:bsd-suc-hstar}
 h_*:=-vf,
 \qquad
 \partial h_*=z_{\mathrm{PT}}-z_K.
\end{equation}
No se introduce una homotopía abstracta como premisa: la cerradura fija
\(u=3cv\), la unidad raíz fija \(a=1\) y el generador \(f\) rellena
directamente la diferencia completa.

\paragraph{Control homotópico equivalente no gobernante.}
\noindent\textbf{Estatuto: CONTROL\_NO\_GOBERNANTE.}
Como comprobación posterior, si se empaqueta el mismo complemento mediante
un lector \(H_{\mathrm{FS}}\), la identidad
\begin{equation}\label{eq:bsd-suc-hfs}
 \partial H_{\mathrm{FS}}+H_{\mathrm{FS}}\partial
 =q_{\mathrm{root}}
\end{equation}
reproduce \eqref{eq:bsd-suc-hstar} al aplicarse a
\(q_{\mathrm{root}}z_K\). Este lector es una notación auxiliar de la
contracción ya realizada por \(h_*=-vf\); no es una obligación anterior ni
puede reabrir la construcción explícita.

\subsection*{Regulador Bockstein derivado de orden arbitrario}

Sea \(K\) un campo de característica cero,
\(\Lambda=K[[T]]\), y sea
\begin{equation}\label{eq:bsd-suc-S}
 S(T):V_{\Lambda}\longrightarrow W_{\Lambda},
 \qquad
 V_{\Lambda}\simeq W_{\Lambda}\simeq\Lambda^{m},
\end{equation}
un diferencial basado que se vuelve invertible sobre \(K((T))\). Se fija
la línea
\begin{equation}\label{eq:bsd-suc-lambda0}
 \lambda_0(S):=(\det_K V)^{-1}\otimes_K\det_K W,
 \quad
 V=V_{\Lambda}/TV_{\Lambda},\quad
 W=W_{\Lambda}/TW_{\Lambda}.
\end{equation}
Una forma de Smith basada tiene la expresión
\begin{equation}\label{eq:bsd-suc-smith}
 U(T)S(T)V(T)
 =\operatorname{diag}
 \bigl(T^{a_1},\ldots,T^{a_s},1,\ldots,1\bigr),
 \qquad a_i\geq1,
\end{equation}
y
\begin{equation}\label{eq:bsd-suc-orden}
 d=\operatorname{ord}_{T}\det S(T)=\sum_{i=1}^{s}a_i.
\end{equation}

La filtración \(T\)-ádica del complejo
\([V_{\Lambda}\xrightarrow{S}W_{\Lambda}]\) produce páginas
\(V_q=E_q^0\), \(W_q=E_q^1\) y Bocksteins
\(\beta_q:V_q\to W_q\). Intrínsecamente,
\begin{equation}\label{eq:bsd-suc-pages}
 A_q:=V_q/V_{q+1},
 \qquad
 B_q:=\ker(W_q\twoheadrightarrow W_{q+1})
      =\operatorname{im}\beta_q,
\end{equation}
y \(\beta_q\) induce
\begin{equation}\label{eq:bsd-suc-beta-q}
 \bar\beta_q:A_q\xrightarrow{\ \sim\ }B_q,
 \qquad
 \tau_q:=\det(\bar\beta_q).
\end{equation}
Las páginas son finitas y verifican
\begin{equation}\label{eq:bsd-suc-carry-orden}
 d=\sum_{q\geq1}q\,\dim A_q
  =\sum_{q\geq1}\dim V_q.
\end{equation}
La primera página ve
\(r_0=\dim V_1\); el acarreo de orden superior es
\begin{equation}\label{eq:bsd-suc-carry-superior}
 c_I^{\deg}=d-r_0=\sum_{q\geq2}\dim V_q.
\end{equation}

La página regular es el isomorfismo inducido
\begin{equation}\label{eq:bsd-suc-pagina-regular}
 \bar S(0):A_0:=V/V_1
 \xrightarrow{\ \sim\ }
 B_0:=\operatorname{im}S(0),
 \qquad
 \tau_{\mathrm{reg}}:=\det\bar S(0).
\end{equation}
Las sucesiones exactas de \eqref{eq:bsd-suc-pages}, con sus signos de
Koszul, pegan estas trivializaciones en un único elemento basado
\begin{equation}\label{eq:bsd-suc-rboc}
 R_{\mathrm{Boc}}^{\mathrm{der}}(S)
 :=\tau_{\mathrm{reg}}\otimes
   \bigotimes_{q\geq1}\tau_q
 \in\lambda_0(S),
\end{equation}
donde el símbolo tensorial abrevia los isomorfismos de pegado de
Knudsen--Mumford.

\begin{theorem}[Identidad determinantal del regulador derivado]
\label{thm:bsd-suc-bockstein}
En la línea basada \(\lambda_0(S)\) se tiene
\begin{equation}\label{eq:bsd-suc-lt}
 R_{\mathrm{Boc}}^{\mathrm{der}}(S)
 =\operatorname{LT}_{T}(S)
 :=\left[T^{-d}\det S(T)\right]_{T=0}.
\end{equation}
La igualdad conserva la página regular, todas las páginas positivas y el
acarreo \eqref{eq:bsd-suc-carry-superior}.
\end{theorem}

\begin{proof}
En las bases Smith de \eqref{eq:bsd-suc-smith}, la página regular y cada
página positiva llevan el generador canónico de su línea. Los signos de
Koszul del pegado son exactamente los necesarios para restituir el orden de
las bases totales, por lo que su producto es el generador positivo de
\(\lambda_0(S)\). Por otra parte,
\[
 \det S(T)=\det U(T)^{-1}\det V(T)^{-1}T^d,
\]
y el término inicial es
\(\det U(0)^{-1}\det V(0)^{-1}\). Al deshacer el cambio de bases, tanto
\eqref{eq:bsd-suc-rboc} como el término inicial se multiplican por ese mismo
factor. Esto prueba \eqref{eq:bsd-suc-lt} sin identificar escalares antes de
fijar la línea y su base.
\end{proof}

Para cada \(q\geq1\), la dualidad de Poitou--Tate construida en la misma
página proporciona
\begin{equation}\label{eq:bsd-suc-jq}
 j_q:B_q\xrightarrow{\ \sim\ }
 A_q^{\vee}\otimes(I^q/I^{q+1}),
\end{equation}
y, tras orientar \(I^q/I^{q+1}\), el emparejamiento
\begin{equation}\label{eq:bsd-suc-altura-derivada}
 h_E^{(q)}(x,y)
 :=\bigl(j_q\bar\beta_q(x)\bigr)(y).
\end{equation}

% Ampliación sustantiva de los comparadores determinantes BSD.
% Módulo destinado al capítulo BSD de la Parte IV.

\section{Comparadores de la unidad determinantal común}
\label{sec:amp-bsd-comparadores}

La igualdad raíz \eqref{eq:bsd-suc-raiz} fija una sola unidad en los canales
Kato y Poitou--Tate. A partir de esa unidad se construyen ahora los
comparadores que transportan el regulador derivado hasta su publicación
escalar. Ningún comparador elige una normalización después de conocer el
término principal: todos proceden de morfismos de complejos, sucesiones
exactas y bases ya orientadas.

\subsection{La cadena de líneas determinantes}
\label{subsec:amp-bsd-lineas}

Para un complejo perfecto basado \(C\), se emplea la convención
\begin{equation}\label{eq:amp-bsd-linea-determinante}
 \lambda(C):=\bigotimes_i
 \det H^i(C)^{(-1)^i}.
\end{equation}
Sean \(\mathcal C_E^K\), \(\mathcal C_E^{\mathrm{Iw}}\),
\(\mathcal C_E^{\mathrm{PT}}\) y \(\mathcal C_E^{\mathrm{loc}}\),
respectivamente, el complejo publicado por el canal de Kato, su modelo
Iwasawa basado, el complejo reducido autodual de Poitou--Tate y el cono de
localización finita--singular. Sus líneas se enlazan mediante
\begin{equation}\label{eq:amp-bsd-cadena-lineas}
 \mathcal L_E^K
 \xrightarrow{\ \mathfrak c_{K/\mathrm{FERS}}\ }
 \mathcal L_E^{\mathrm{Iw}}
 \xrightarrow{\ \mathfrak c_{\mathrm{PT}}\ }
 \mathcal L_E^{\mathrm{ht}}
 \xrightarrow{\ \mathfrak c_{\mathrm{STS}}\ }
 \mathcal L_E^{\mathrm{fin}}
 \xrightarrow{\ \mathfrak c_{\mathrm{tors}}\ }
 \mathcal L_E^{\mathrm{ar}}
 \xrightarrow{\ \mathfrak c_{\infty}\ }
 \mathcal L_E.
\end{equation}
Las siglas \(\mathrm{STS}\) designan en esta fórmula el paso conjunto de
Smith, Tamagawa y Tate--Shafarevich. Todas las flechas de
\eqref{eq:amp-bsd-cadena-lineas} son isomorfismos de líneas orientadas.

\subsection{Comparación de cadenas Kato--FERS desde la unidad común}
\label{subsec:amp-bsd-comparacion-cadenas-kato-fers}

La equivalencia que origina el primer comparador puede construirse antes de
tomar determinantes. Sea \(F^q\mathcal C\) la filtración por profundidad
nonádica, identificada con la filtración \(T\)-ádica mediante el acarreo, y
escríbase
\begin{equation}\label{eq:amp-bsd-complejos-kato-iwasawa}
 \mathcal C_E^K
 :=\operatorname{Tot}\!\left(
 P_E^K\widehat\Pi_0G_{m,n}\right),
 \qquad
 \mathcal C_E^{\mathrm{Iw}}
 :=[V_\Lambda\xrightarrow{S_E(T)}W_\Lambda].
\end{equation}
Sea
\(\operatorname{car}_T:G_{m,n}\to\mathcal C_E^{\mathrm{Iw}}\) el lector
de acarreo que asigna a una ruta \(\gamma\) el generador determinado por sus
caras inicial y final, multiplicado por
\(T^{\operatorname{Carr}(\gamma)}\). La multiplicación de concatenaciones en
el modelo Iwasawa se denota por \(\mu_{\mathrm{Iw}}\). La diagonal
Alexander--Whitney y la publicación \(P_E^K\) definen entonces, antes de
tomar cohomología o determinantes, el mapa de cadenas
\begin{equation}\label{eq:amp-bsd-phi-kato-fers}
\begin{aligned}
 \Phi_{K/\mathrm{FERS}}&:
 \mathcal C_E^K\longrightarrow\mathcal C_E^{\mathrm{Iw}},\\
 \Phi_{K/\mathrm{FERS}}
 &:=
 \mu_{\mathrm{Iw}}\circ
 \bigl(P_E^K\widehat\otimes\operatorname{car}_T\bigr)
 \circ\Delta_{\mathrm{AW}},\\
 \Phi_{K/\mathrm{FERS}}(z_K)&=z_{\mathrm{Iw}}.
\end{aligned}
\end{equation}
donde \(z_{\mathrm{Iw}}\) es el generador raíz inducido por \(1_E\). Así,
cada simplex normalizado se envía a la suma de sus cortes
Alexander--Whitney, conservando en cada sumando las dos caras y la potencia
de acarreo. La fórmula no consulta el término principal analítico.

\begin{lemma}[Criterio generador a generador para la equivalencia]
\label{lem:amp-bsd-bases-emparejadas-kato-fers}
En cada grado \(i\), sea \(\mathcal B_i^K\) la base finita de simplices
normalizados, ordenada por profundidad, caras y acarreo, y sea
\(\mathcal B_i^{\mathrm{Iw}}\) la base formada por los generadores con las
mismas caras inicial y final en el modelo Iwasawa. Para una aplicación
candidata
\[
 \upsilon_i:\mathcal B_i^K\longrightarrow
 \mathcal B_i^{\mathrm{Iw}},
\]
defínase el residual de grado cero
\begin{equation}\label{eq:amp-bsd-residual-generador-graduado}
 \varepsilon_i(b)
 :=\Phi_{K/\mathrm{FERS}}^i(b)\bmod T-\upsilon_i(b).
\end{equation}
Si \(\upsilon_i\) es una biyección y
\(\varepsilon_i(b)=0\) para todo \(b\in\mathcal B_i^K\), entonces los dos
módulos completados del grado \(i\) son libres del mismo rango
\(d_i=|\mathcal B_i^K|\) y, para cada generador,
\begin{equation}\label{eq:amp-bsd-phi-en-cada-generador}
 \Phi_{K/\mathrm{FERS}}^i(b)
 =\upsilon_i(b)
  +T\sum_{c\in\mathcal B_i^{\mathrm{Iw}}}
    \lambda_{c,b}(T)c,
 \qquad \lambda_{c,b}(T)\in\Lambda.
\end{equation}
En particular,
\(\operatorname{gr}^{F}\Phi_{K/\mathrm{FERS}}^i=I\) sobre todos los
generadores, no sólo sobre la unidad raíz.
\end{lemma}

\begin{proof}
La anulación de \eqref{eq:amp-bsd-residual-generador-graduado} dice que la
columna de \(b\) tiene término constante \(\upsilon_i(b)\). Los restantes
coeficientes pertenecen a \(T\Lambda\), lo que da
\eqref{eq:amp-bsd-phi-en-cada-generador}. La biyectividad identifica las
bases y fija el mismo rango. En esas bases, la reducción módulo \(T\) de la
matriz de \(\Phi^i\) es la identidad; de ahí
\(\operatorname{gr}^{F}\Phi^i=I\).
\end{proof}

La igualdad raíz
\(\Phi_{K/\mathrm{FERS}}(z_K)=z_{\mathrm{Iw}}\) fija la normalización común.
La escritura de \(P_E^K\) sobre cada simplex normalizado, la biyección
\(\upsilon_i\) y los residuales \(\varepsilon_i(b)\) proporcionan un
control generador a generador de la equivalencia ya construida; no eligen ni
gobiernan la conclusión BSD.

\begin{lemma}[Equivalencia filtrada y normalización de la raíz]
\label{lem:amp-bsd-equivalencia-kato-fers}
Para el mapa de cadenas construido en
\eqref{eq:amp-bsd-phi-kato-fers}, la aplicación \(\upsilon_i\) es la
correspondencia de bases inducida por las caras y el acarreo de cada
simplejo. La fórmula de Alexander--Whitney verifica en cada grado el
criterio del lema \ref{lem:amp-bsd-bases-emparejadas-kato-fers}; éste es un
control de la equivalencia ya construida y no una premisa adicional.
La aplicación \(\Phi_{K/\mathrm{FERS}}\) es una equivalencia basada de
complejos perfectos. En las bases ordenadas por profundidad, sus matrices en
cada grado tienen la forma
\begin{equation}\label{eq:amp-bsd-matriz-kato-fers}
 [\Phi_{K/\mathrm{FERS}}^i](T)=I+TB_i(T),
 \qquad B_i(T)\in M_{d_i}(\Lambda).
\end{equation}
Por tanto, la unidad
\begin{equation}\label{eq:amp-bsd-rho-determinantal}
 \rho_{E,\ell}(T)
 :=\prod_i\det\!\left(I+TB_i(T)\right)^{(-1)^i}
 \in\Lambda^\times
\end{equation}
satisface \(\rho_{E,\ell}(0)=1\).
\end{lemma}

\begin{proof}
La identidad simplicial
\(\partial\Delta_{\mathrm{AW}}
=(\partial\otimes I+(-1)^{\deg}I\otimes\partial)
\Delta_{\mathrm{AW}}\), junto con la compatibilidad de
\(P_E^K\) con caras y degeneraciones, demuestra que
\eqref{eq:amp-bsd-phi-kato-fers} conmuta con los diferenciales. En el graduado
asociado, el lema
\ref{lem:amp-bsd-bases-emparejadas-kato-fers} identifica todos los
generadores y da la identidad. En esas bases,
\eqref{eq:amp-bsd-phi-en-cada-generador} es precisamente
\eqref{eq:amp-bsd-matriz-kato-fers}.

La completitud de \(\Lambda=K[[T]]\) hace converger la serie
\begin{equation}\label{eq:amp-bsd-inversa-kato-fers}
 (I+TB_i(T))^{-1}
 =\sum_{n\geq0}(-TB_i(T))^n;
\end{equation}
estas inversas también conmutan con los diferenciales, por lo que forman la
inversa basada de \(\Phi_{K/\mathrm{FERS}}\). Al tomar el determinante
alternado se obtiene \eqref{eq:amp-bsd-rho-determinantal}; evaluar en
\(T=0\) da uno. Así, la normalización no se postula a partir del término
principal: es la coordenada constante de una equivalencia que es la identidad
sobre la unidad raíz.
\end{proof}

El primer comparador es el determinante de
\(\Phi_{K/\mathrm{FERS}}\). Si \(\mathfrak z_K\) denota el elemento
procedente de la unidad \(1_E\), su imagen se caracteriza por
\begin{equation}\label{eq:amp-bsd-comparador-kato}
 \mathfrak c_{K/\mathrm{FERS}}(\mathfrak z_K)
 =\left[T^{-d}\rho_{E,\ell}(T)\det S_E(T)\right]_{T=0},
 \qquad
 \rho_{E,\ell}(0)=1.
\end{equation}
La unidad de \(\rho_{E,\ell}\) es orientada y vale uno en la raíz; por ello no
altera el orden ni la base. El teorema
\ref{thm:bsd-suc-bockstein} transforma \eqref{eq:amp-bsd-comparador-kato} en
\begin{equation}\label{eq:amp-bsd-kato-bockstein}
 \mathfrak c_{K/\mathrm{FERS}}(\mathfrak z_K)
 =R_{\mathrm{Boc}}^{\mathrm{der}}(S_E)
 =\tau_{\mathrm{reg}}\otimes\bigotimes_{q\geq1}\tau_q.
\end{equation}

El comparador de Poitou--Tate se construye página a página. Para
\(q\geq1\), el isomorfismo \(j_q\) de
\eqref{eq:bsd-suc-jq} y el Bockstein reducido
\(\bar\beta_q\) inducen
\begin{equation}\label{eq:amp-bsd-comparador-pt-q}
 \mathfrak c_{\mathrm{PT},q}(\tau_q)
 :=\det\!\left(j_q\circ\bar\beta_q\right)
 =\det h_E^{(q)}.
\end{equation}
La dirección regular se transporta mediante
\(\det\bar S_E(0)=\tau_{\mathrm{reg}}\). El pegado de Knudsen--Mumford y los
signos de Koszul definidos en \eqref{eq:bsd-suc-rboc} proporcionan el único
isomorfismo orientado
\begin{equation}\label{eq:amp-bsd-comparador-pt}
 \mathfrak c_{\mathrm{PT}}
 \left(\tau_{\mathrm{reg}}\otimes\bigotimes_{q\geq1}\tau_q\right)
 =\tau_{\mathrm{reg}}^{\mathrm{ht}}
  \otimes\bigotimes_{q\geq1}\det h_E^{(q)}.
\end{equation}
En la primera página estable, la forma \(h_E^{(1)}\) es la forma
Poincaré--Néron--Tate sobre la red libre de Mordell--Weil; su determinante es
\begin{equation}\label{eq:amp-bsd-regulador-nt}
 \det h_E^{(1)}=\operatorname{Reg}(E)
\end{equation}
en la base inducida por la unidad raíz.

\begin{lemma}[Despliegue Poitou--Tate del grado aritmético]
\label{lem:amp-bsd-grado-pt}
Sea
\[
 \operatorname{Flag}_{\mathrm{Boc}}(E)
 :=\bigoplus_{q\geq1}V_q
\]
la bandera finita de páginas positivas, cada una con la base heredada de la
unidad raíz. La reducción ortogonal y los isomorfismos \(j_q\) construyen un
isomorfismo basado
\begin{equation}\label{eq:amp-bsd-psi-pt}
 \Psi_{\mathrm{PT}}:
 \bigl(E(\mathbb Q)/E(\mathbb Q)_{\mathrm{tors}}\bigr)\otimes K
 \xrightarrow{\ \sim\ }\operatorname{Flag}_{\mathrm{Boc}}(E).
\end{equation}
En particular,
\begin{equation}\label{eq:amp-bsd-rango-bandera}
 \operatorname{rank}E(\mathbb Q)
 =\sum_{q\geq1}\dim_KV_q
 =\operatorname{ord}_T\det S_E(T).
\end{equation}
\end{lemma}

\begin{proof}
Se filtra el complejo reducido por las potencias del ideal de aumento. En el
cociente de nivel \(q\), la dirección regular ya ha sido separada por
\(\bar S_E(0)\), y el único bloque superviviente es \(V_q\). La equivalencia
\(X^{\mathrm{orth}}\) identifica el bloque opuesto con su dual, mientras
\(j_q\bar\beta_q\) proporciona el emparejamiento perfecto que levanta la
base del cociente al nivel anterior. Comenzando en la última página no nula y
repitiendo este levantamiento, se obtiene \(\Psi_{\mathrm{PT}}\).

En las bases ordenadas por profundidad, la matriz del mapa resultante es
triangular por bloques y todos sus bloques diagonales son identidades: en el
primer bloque porque \(z_K\) y \(z_{\mathrm{PT}}\) tienen la misma raíz, y
en los restantes porque \(j_q\) y \(\bar\beta_q\) son isomorfismos basados.
Por ello \(\Psi_{\mathrm{PT}}\) es invertible y no pierde ninguna página.
Tomando dimensiones y usando
\eqref{eq:bsd-suc-carry-orden} se obtiene
\eqref{eq:amp-bsd-rango-bandera}. Al tomar determinantes, el mismo
levantamiento es exactamente el pegado de
\eqref{eq:amp-bsd-comparador-pt}; así, la igualdad de rangos y la igualdad de
líneas proceden de un solo mapa basado.
\end{proof}

\subsection{Triángulo de localización, rango de Smith y factores finitos}
\label{subsec:amp-bsd-smith-sha}

Para cada lugar \(v\), la condición local finita es un subcomplejo basado
\(\mathcal C_{E,v}^{\mathrm{fin}}\subset\mathcal C_{E,v}\), y se define
\(\mathcal C_{E,v}^{\mathrm{sing}}\) por el triángulo
\begin{equation}\label{eq:amp-bsd-triangulo-local}
 \mathcal C_{E,v}^{\mathrm{fin}}
 \xrightarrow{\ i_v\ }\mathcal C_{E,v}
 \longrightarrow\mathcal C_{E,v}^{\mathrm{sing}}
 \longrightarrow\mathcal C_{E,v}^{\mathrm{fin}}[1].
\end{equation}
Si \(\operatorname{res}_v\) denota la restricción global, el mapa de
localización es, a nivel de cadenas,
\begin{equation}\label{eq:amp-bsd-mapa-localizacion}
 \ell_E:
 \mathcal C_E^{\mathrm{glob,red}}\oplus
 \bigoplus_v\mathcal C_{E,v}^{\mathrm{fin}}
 \longrightarrow\bigoplus_v\mathcal C_{E,v},
 \qquad
 \ell_E\bigl(x,(y_v)_v\bigr)
 :=\bigl(\operatorname{res}_v x-i_vy_v\bigr)_v .
\end{equation}
La definición por cono del complejo Selmer reducido da el triángulo
\begin{equation}\label{eq:amp-bsd-triangulo-localizacion}
 \mathcal C_E^{\mathrm{red}}
 \longrightarrow
 \mathcal C_E^{\mathrm{glob,red}}\oplus
 \bigoplus_v\mathcal C_{E,v}^{\mathrm{fin}}
 \xrightarrow{\ \ell_E\ }
 \bigoplus_v\mathcal C_{E,v}
 \longrightarrow\mathcal C_E^{\mathrm{red}}[1].
\end{equation}
Por tanto, la exactitud local--global procede del cono de un mapa escrito en
\eqref{eq:amp-bsd-mapa-localizacion}; no se añade después como una relación
entre cardinales.

Se cancela en \eqref{eq:amp-bsd-triangulo-localizacion} el sumando raíz
común, que es el mismo en ambos canales por
\eqref{eq:bsd-suc-raiz}. El complemento resultante se denota por
\begin{equation}\label{eq:amp-bsd-cono-local-reducido}
 \mathcal C_E^{\mathrm{loc,red}}
 :=q_{\mathrm{root}}\operatorname{Cone}(\ell_E)[-1].
\end{equation}
Tras separar la red libre de Mordell--Weil y su dual mediante
\(X^{\mathrm{orth}}\), una elección de las bases integrales ya orientadas
representa este complejo por
\begin{equation}\label{eq:amp-bsd-complejo-smith}
 \mathcal C_E^{\mathrm{loc,red}}
 \simeq[M_E\xrightarrow{\ D_E\ }N_E].
\end{equation}

\begin{lemma}[Aciclicidad racional y rango completo]
\label{lem:amp-bsd-aciclicidad-rango}
El complejo de \eqref{eq:amp-bsd-complejo-smith} es acíclico después de
tensorizar con \(\mathbb Q\). En particular,
\[
 D_E\otimes\mathbb Q:M_E\otimes\mathbb Q
 \xrightarrow{\ \sim\ }N_E\otimes\mathbb Q
\]
es un isomorfismo; \(M_E\) y \(N_E\) tienen el mismo rango y \(D_E\) tiene
rango completo.
\end{lemma}

\begin{proof}
Sea \(x\) un ciclo del complejo
\(\mathcal C_E^{\mathrm{loc,red}}\otimes\mathbb Q\). La reducción ha
eliminado el sumando \(p_{\mathrm{root}}\), luego
\(q_{\mathrm{root}}x=x\). En las bases integrales orientadas empleadas en
\eqref{eq:amp-bsd-complejo-smith}, el complemento finita--singular se
descompone en los bloques normales del
\cref{lem:bsd-suc-forma-normal}. En cada bloque \(\lambda\), el coeficiente
raíz ya es cero y el ciclo se escribe
\[
 x_{\lambda}=u_{\lambda}e_{\lambda}
                  +v_{\lambda}b_{\lambda},
 \qquad
 \partial x_{\lambda}
   =(u_{\lambda}-3c_{\lambda}v_{\lambda})g_{\lambda}=0.
\]
Como \(g_{\lambda}\) pertenece a la base libre, se tiene
\(u_{\lambda}=3c_{\lambda}v_{\lambda}\), y la fórmula explícita del
diferencial da
\[
 x_{\lambda}
   =v_{\lambda}(3c_{\lambda}e_{\lambda}+b_{\lambda})
   =\partial(v_{\lambda}f_{\lambda}).
\]
La suma es finita en cada grado, de modo que
\(x=\partial\bigl(\sum_{\lambda}v_{\lambda}f_{\lambda}\bigr)\).
Todo ciclo es, por tanto, una frontera por la misma forma normal que produjo
el filler \eqref{eq:bsd-suc-hstar}, sin introducir
\(H_{\mathrm{FS}}\) como premisa. La igualdad
\(p_{\mathrm{root}}z_K=z_{\mathrm{PT}}\) es la que permite cancelar la
dirección raíz antes de este cálculo; sin ella quedaría una clase de grado
cero. La equivalencia autodual
\(X^{\mathrm{orth}}:\mathcal C_E^{\mathrm{red}}\simeq
D_3(\mathcal C_E^{\mathrm{red}})\) empareja los dos grados restantes y
transporta sus bases con orientación opuesta complementaria. De aquí se
obtiene \(\operatorname{rank}M_E=\operatorname{rank}N_E\).
Finalmente, la aciclicidad de un complejo de dos términos de igual rango
equivale a que \(D_E\otimes\mathbb Q\) sea invertible.
\end{proof}

El lema permite aplicar la forma normal de Smith sin presuponer ninguno de
sus divisores. Existen cambios de base orientados \(U,V\) tales que
\begin{equation}\label{eq:amp-bsd-morfismo-smith}
 U D_E V=\operatorname{diag}(s_1,\ldots,s_t),
 \qquad s_i\in\mathbb Z\setminus\{0\}.
\end{equation}
Por tanto,
\begin{equation}\label{eq:amp-bsd-grupo-smith}
 F_E:=\operatorname{coker}D_E
 \simeq\bigoplus_{i=1}^{t}\mathbb Z/|s_i|\mathbb Z,
 \qquad
 |F_E|=\prod_{i=1}^{t}|s_i|<\infty.
\end{equation}

\begin{lemma}[Sucesión finita de localización]
\label{lem:amp-bsd-exactitud-sha-tamagawa}
El fragmento torsional de la sucesión larga de cohomología de
\eqref{eq:amp-bsd-triangulo-localizacion} es
\begin{equation}\label{eq:amp-bsd-exacta-sha-tamagawa}
 0\longrightarrow\operatorname{Sha}(E)
 \xrightarrow{\ \iota_{\mathrm{Sha}}\ }F_E
 \xrightarrow{\ \partial_{\mathrm{loc}}\ }
 \bigoplus_v\Phi_v(\mathbb F_v)
 \longrightarrow0,
 \qquad c_v=|\Phi_v(\mathbb F_v)|.
\end{equation}
\end{lemma}

\begin{proof}
Un cociclo global representa una clase de
\(\operatorname{Sha}(E)\) precisamente cuando su restricción en cada lugar
admite una primitiva en
\(\mathcal C_{E,v}^{\mathrm{fin}}\). La clase del par formado por ese
cociclo y sus primitivas locales en el cono de
\eqref{eq:amp-bsd-mapa-localizacion} define
\(\iota_{\mathrm{Sha}}\). Si esta clase es una frontera, la componente global es una
frontera y la clase de Tate--Shafarevich es nula; así,
\(\iota_{\mathrm{Sha}}\) es inyectiva.

El mapa \(\partial_{\mathrm{loc}}\) toma una clase del cono, la restringe a
cada cociente
\(\mathcal C_{E,v}^{\mathrm{sing}}\) de
\eqref{eq:amp-bsd-triangulo-local} y conserva su clase de componente. La
identificación
\[
 H^0(\mathcal C_{E,v}^{\mathrm{sing}})_{\mathrm{tors}}
 \simeq\Phi_v(\mathbb F_v)
\]
es inducida por el cociente finita--singular, no por el orden del grupo.
Una clase del cono tiene residuo cero exactamente cuando sus representantes
locales pueden elegirse en los subcomplejos finitos; por la definición
anterior, ese núcleo es la imagen de \(\iota_{\mathrm{Sha}}\).
Finalmente, la última flecha de
\eqref{eq:amp-bsd-triangulo-local} levanta cada clase de componente al
complejo local. La exactitud del cono
\eqref{eq:amp-bsd-triangulo-localizacion}, después de retirar los dos
sumandos libres emparejados por \(X^{\mathrm{orth}}\), convierte ese
levantamiento en una preimagen por \(\partial_{\mathrm{loc}}\). Esto prueba
la sobreyectividad y, por consiguiente, toda
\eqref{eq:amp-bsd-exacta-sha-tamagawa}.
\end{proof}

Ahora, y sólo después del lema
\ref{lem:amp-bsd-aciclicidad-rango}, el grupo \(F_E\) es finito. La
inyección de \eqref{eq:amp-bsd-exacta-sha-tamagawa} demuestra entonces la
finitud de \(\operatorname{Sha}(E)\). Al tomar órdenes en esa sucesión
exacta se obtiene
\begin{equation}\label{eq:amp-bsd-cardinal-smith}
 |F_E|=|\operatorname{Sha}(E)|\prod_v c_v.
\end{equation}
Éste es el comparador
\(\mathfrak c_{\mathrm{STS}}\): transporta los factores elementales de Smith
a los factores Tate--Shafarevich y Tamagawa conservando su orden en la línea,
en lugar de insertar sus cardinales una vez calculado el término principal.

La parte torsional se construye aplicando el determinante a la red de
Mordell--Weil y a su dual de Pontryagin. Los dos factores finitos son duales y
aportan
\begin{equation}\label{eq:amp-bsd-comparador-torsion}
 \mathfrak c_{\mathrm{tors}}:quad
 u\longmapsto
 \frac{u}{|E(\mathbb Q)_{\mathrm{tors}}|^2}.
\end{equation}
El comparador arquimediano se obtiene integrando la diferencial de Néron
orientada sobre el ciclo real basado:
\begin{equation}\label{eq:amp-bsd-comparador-periodo}
 \Omega_E:=\int_{E(\mathbb R)}|\omega_E|,
 \qquad
 \mathfrak c_\infty:quad u\longmapsto\frac{u}{\Omega_E}.
\end{equation}
La diferencial y el ciclo se fijan antes de evaluar \(L(E,s)\); por ello
\eqref{eq:amp-bsd-comparador-periodo} no normaliza retrospectivamente la
sección analítica.

\subsection{Separación causal de primitivas, identidades y consecuencias}
\label{subsec:amp-bsd-separacion-causal}

El criterio generador a generador del
lema \ref{lem:amp-bsd-bases-emparejadas-kato-fers} audita en cada grado la
equivalencia ya construida. No condiciona su existencia. El orden lógico
del paquete es
\begin{equation}\label{eq:amp-bsd-diagrama-causal}
\begin{array}{c}
 \begin{gathered}
  1_E,\ \Delta_{\mathrm{AW}},\ P_E^K,\ P_E^{\mathrm{PT}},
  \ h_*=-vf,\ X^{\mathrm{orth}},\\
  \{\,\mathcal C_{E,v}^{\mathrm{fin}}\to\mathcal C_{E,v}\,\}_v,\\
  \text{bases de Mordell--Weil y Pontryagin;}\\
  \text{diferencial y ciclo de Néron}
 \end{gathered}
 \\[2mm]\Downarrow\\[-1mm]
 \begin{gathered}
  \Phi_{K/\mathrm{FERS}}\text{ equivalencia},\quad
  \rho_{E,\ell}(0)=1,\quad
  p_{\mathrm{root}}z_K=z_{\mathrm{PT}},\\
  R_{\mathrm{Boc}}^{\mathrm{der}}=\operatorname{LT}_T(S_E),\quad
  \det(j_q\bar\beta_q)=\det h_E^{(q)},\\
  H^*(\mathcal C_E^{\mathrm{loc,red}}\otimes\mathbb Q)=0,\quad
  D_E\otimes\mathbb Q\text{ isomorfismo},\\
  0\to\operatorname{Sha}(E)\to F_E
  \to\bigoplus_v\Phi_v(\mathbb F_v)\to0
 \end{gathered}
 \\[2mm]\Downarrow\\[-1mm]
 \begin{gathered}
  d_{\mathrm{an}}=\operatorname{rank}E(\mathbb Q),\qquad
  \operatorname{Reg}(E),\qquad
  |\operatorname{Sha}(E)|\prod_vc_v,\\
  |E(\mathbb Q)_{\mathrm{tors}}|^{-2},\qquad
  \Omega_E^{-1},\qquad
  \text{igualdad del término principal}.
 \end{gathered}
\end{array}
\end{equation}
La primera fila contiene únicamente objetos y mapas ya basados, junto con
la verificación generador a generador indicada; la segunda, identidades
demostradas por esos mapas; la tercera, sus consecuencias sobre la línea
determinante. En particular, ningún factor de la última fila se
emplea para normalizar retroactivamente una flecha de la primera.

\begin{proposition}[Paquete exacto de comparadores]
\label{prop:amp-bsd-paquete-comparadores-exacto}
Los comparadores de
\eqref{eq:amp-bsd-cadena-lineas} proceden de un único paquete de mapas de
complejos y satisfacen, en el orden de construcción,
\begin{equation}\label{eq:amp-bsd-paquete-identidades}
 \begin{gathered}
 \Phi_{K/\mathrm{FERS}}(z_K)=z_{\mathrm{Iw}},\qquad
 \operatorname{gr}^{F}\Phi_{K/\mathrm{FERS}}=I,\\
 \det\nolimits_{\mathrm{alt}}\Phi_{K/\mathrm{FERS}}
   =\rho_{E,\ell}(T),\qquad \rho_{E,\ell}(0)=1,\\
 \det(j_q\bar\beta_q)=\det h_E^{(q)}\quad(q\geq1),\\
 \Psi_{\mathrm{PT}}:
  \bigl(E(\mathbb Q)/E(\mathbb Q)_{\mathrm{tors}}\bigr)\otimes K
  \xrightarrow{\sim}\operatorname{Flag}_{\mathrm{Boc}}(E),\\
 p_{\mathrm{root}}z_K=z_{\mathrm{PT}},\qquad
 h_*=-vf,\quad \partial h_*=z_{\mathrm{PT}}-z_K,\\
 H^*(\mathcal C_E^{\mathrm{loc,red}}\otimes\mathbb Q)=0,\qquad
 U D_E V=\operatorname{diag}(s_1,\ldots,s_t),\\
 0\longrightarrow\operatorname{Sha}(E)
  \longrightarrow F_E
  \longrightarrow\displaystyle\bigoplus_v\Phi_v(\mathbb F_v)
  \longrightarrow0 .
 \end{gathered}
\end{equation}
La composición inducida en líneas determinantes es exactamente
\begin{equation}\label{eq:amp-bsd-paquete-composicion}
 \det(\Phi_{K/\mathrm{FERS}})
 \xrightarrow{\ \det(j_q\bar\beta_q)\ }
 \det(D_E)
 \xrightarrow{\ \mathrm{Smith}\ }
 \frac{|\operatorname{Sha}(E)|\prod_vc_v}
      {|E(\mathbb Q)_{\mathrm{tors}}|^2\,\Omega_E},
\end{equation}
con los signos de Knudsen--Mumford y las orientaciones de
\eqref{eq:amp-bsd-cadena-lineas}. En particular, el miembro derecho de
\eqref{eq:amp-bsd-paquete-composicion} es la coordenada final del elemento
transportado y no un dato con el que se construya alguna de las flechas.
\end{proposition}

\begin{proof}
La primera fila de \eqref{eq:amp-bsd-paquete-identidades} es
\eqref{eq:amp-bsd-phi-kato-fers} y
\eqref{eq:amp-bsd-matriz-kato-fers}--
\eqref{eq:amp-bsd-rho-determinantal}. La segunda es
\eqref{eq:amp-bsd-comparador-pt-q} y el lema
\ref{lem:amp-bsd-grado-pt}. La tercera combina la igualdad de la raíz con
el filler explícito \eqref{eq:bsd-suc-hstar}; su extensión lineal en la
forma normal induce el lector auxiliar utilizado en el lema
\ref{lem:amp-bsd-aciclicidad-rango}. Al restringir al complemento de la
raíz, esa reducción contrae todo ciclo racional. La última fila es
\eqref{eq:amp-bsd-morfismo-smith} y
\eqref{eq:amp-bsd-exacta-sha-tamagawa}. Aplicar el funtor determinante a
esas identidades, en ese mismo orden, da
\eqref{eq:amp-bsd-paquete-composicion}. Ninguna igualdad se obtiene
comparando primero los dos escalares finales.
\end{proof}

\begin{lemma}[Independencia causal del término principal]
\label{lem:amp-bsd-independencia-termino-principal}
Sea \(\mathscr P_E\) el conjunto de datos
\begin{equation}\label{eq:amp-bsd-primitivas-paquete}
 \mathscr P_E:=
 \left(
  1_E,\Delta_{\mathrm{AW}},P_E^K,P_E^{\mathrm{PT}},
  h_*=-vf,X^{\mathrm{orth}},
  \{\mathcal C_{E,v}^{\mathrm{fin}}\xrightarrow{i_v}
     \mathcal C_{E,v}\}_v,
  \text{bases y orientaciones},
  \omega_E,\gamma_E^{\mathbb R}
 \right),
\end{equation}
donde \(\omega_E\) es la diferencial de Néron y
\(\gamma_E^{\mathbb R}\) el ciclo real orientado. Todos los mapas del
paquete de la proposición
\ref{prop:amp-bsd-paquete-comparadores-exacto} se determinan por
\(\mathscr P_E\). No forman parte de \(\mathscr P_E\) los numerales
\[
 L^{(r)}(E,1),\quad r,\quad\operatorname{Reg}(E),\quad
 |\operatorname{Sha}(E)|,\quad(c_v)_v,\quad
 |E(\mathbb Q)_{\mathrm{tors}}|,\quad\Omega_E.
\]
Esos valores se obtienen, respectivamente, al evaluar la sección analítica,
tomar el grado de la bandera, calcular determinantes y formas de Smith,
tomar órdenes de grupos finitos e integrar
\(\omega_E\) sobre \(\gamma_E^{\mathbb R}\).
\end{lemma}

\begin{proof}
La fórmula \eqref{eq:amp-bsd-phi-kato-fers} usa sólo
\(1_E,\Delta_{\mathrm{AW}},P_E^K\) y el acarreo. Los mapas Bockstein y
Poitou--Tate usan la filtración, \(P_E^{\mathrm{PT}}\), las bases y la
autodualidad. El cono de \eqref{eq:amp-bsd-mapa-localizacion} usa los
mapas \(i_v\) y las restricciones locales; su contracción y su forma de
Smith usan el filler explícito \eqref{eq:bsd-suc-hstar}, su extensión
lineal auxiliar, \(X^{\mathrm{orth}}\) y las bases integrales. Los dos
últimos comparadores usan las redes torsionales y el
par \((\omega_E,\gamma_E^{\mathbb R})\). Ésta es una enumeración exhaustiva
de los argumentos que aparecen en las fórmulas constructoras. Los
numerales listados en el enunciado aparecen sólo después de aplicar
cohomología, determinante, cardinal o integración, de modo que no pueden
seleccionar retrospectivamente los mapas.
\end{proof}

\subsection{Composición, grados y término principal}
\label{subsec:amp-bsd-composicion}

Defínase el comparador total por la composición en el orden escrito:
\begin{equation}\label{eq:amp-bsd-comparador-total}
 \mathfrak C_E:=
 \mathfrak c_\infty\circ
 \mathfrak c_{\mathrm{tors}}\circ
 \mathfrak c_{\mathrm{STS}}\circ
 \mathfrak c_{\mathrm{PT}}\circ
 \mathfrak c_{K/\mathrm{FERS}}.
\end{equation}
Como \(p_{\mathrm{root}}z_K=z_{\mathrm{PT}}\), todas las flechas reciben la
misma unidad basada. El filler \(h_*=-vf\) de
\eqref{eq:bsd-suc-hstar} rellena explícitamente la diferencia; su extensión
lineal en la base normal produce el lector homotópico auxiliar y prueba que
ninguna clase adicional modifica el elemento de la línea durante el
transporte. El lector no es una premisa separada.

\begin{theorem}[Construcción de los comparadores basados]
\label{thm:amp-bsd-comparadores-construidos}
Para toda curva elíptica \(E/\mathbb Q\) con los complejos, bases,
orientaciones y datos locales definidos anteriormente, las aplicaciones
\eqref{eq:amp-bsd-phi-kato-fers},
\eqref{eq:amp-bsd-comparador-pt},
\eqref{eq:amp-bsd-triangulo-localizacion},
\eqref{eq:amp-bsd-comparador-torsion} y
\eqref{eq:amp-bsd-comparador-periodo}, junto con los lemas
\ref{lem:amp-bsd-aciclicidad-rango} y
\ref{lem:amp-bsd-exactitud-sha-tamagawa}, construyen los cinco comparadores de
\eqref{eq:amp-bsd-cadena-lineas}. Su composición preserva la unidad, el grado
y la orientación, y satisface
\begin{equation}\label{eq:amp-bsd-igualdad-grados}
 d_{\mathrm{an}}=\operatorname{ord}_T\det S_E(T)
 =d_{\mathrm{ar}}=\operatorname{rank}E(\mathbb Q).
\end{equation}
Además, \(\operatorname{Sha}(E)\) es finito y la igualdad del elemento
distinguido en \(\mathcal L_E\) se publica como
\begin{equation}\label{eq:amp-bsd-termino-principal}
 \frac{L^{(r)}(E,1)}{r!\,\Omega_E}
 =\frac{|\operatorname{Sha}(E)|\,\operatorname{Reg}(E)
        \prod_v c_v}
       {|E(\mathbb Q)_{\mathrm{tors}}|^2},
 \qquad r=\operatorname{rank}E(\mathbb Q).
\end{equation}
\end{theorem}

\begin{proof}
El comparador Kato--FERS y la condición
\(\rho_{E,\ell}(0)=1\) identifican el orden analítico con
\(d=\operatorname{ord}_T\det S_E(T)\) y el término inicial con
\(R_{\mathrm{Boc}}^{\mathrm{der}}(S_E)\). La identidad
\eqref{eq:bsd-suc-lt} conserva la página regular y cada página positiva. Los
isomorfismos \(j_q\) transportan esas páginas a las alturas derivadas. El
lema \ref{lem:amp-bsd-grado-pt} construye el isomorfismo de la bandera con la
red libre de Mordell--Weil y, junto con
\eqref{eq:bsd-suc-carry-orden}, prueba
\eqref{eq:amp-bsd-igualdad-grados}.

El lema \ref{lem:amp-bsd-aciclicidad-rango} deduce la aciclicidad racional
directamente del filler explícito \eqref{eq:bsd-suc-hstar}, de su reducción
lineal del complemento, de la cancelación de la raíz común y de
\(X^{\mathrm{orth}}\); de ahí se obtiene el rango completo de \(D_E\), no
como hipótesis adicional. La forma de Smith
\eqref{eq:amp-bsd-morfismo-smith} produce entonces el grupo finito \(F_E\).
El lema \ref{lem:amp-bsd-exactitud-sha-tamagawa} extrae de los triángulos
\eqref{eq:amp-bsd-triangulo-local} y
\eqref{eq:amp-bsd-triangulo-localizacion} la sucesión exacta finita; sólo
después se deducen la finitud de \(\operatorname{Sha}(E)\) y
\eqref{eq:amp-bsd-cardinal-smith}. Finalmente,
\eqref{eq:amp-bsd-regulador-nt},
\eqref{eq:amp-bsd-comparador-torsion} y
\eqref{eq:amp-bsd-comparador-periodo} publican, respectivamente, el regulador,
el cociente torsional y el período. La igualdad raíz elimina cualquier unidad
relativa entre las dos publicaciones. Al evaluar el mismo elemento basado en
ambos extremos de \eqref{eq:amp-bsd-comparador-total}, se obtiene
\eqref{eq:amp-bsd-termino-principal}.
\end{proof}

\subsection{Premisas y criterios de refutación}
\label{subsec:amp-bsd-falsadores}

La construcción parte de la unidad genealógica \(1_E\) y de la diagonal de
Alexander--Whitney. Emplea además las
publicaciones acopladas, todos los
Bocksteins y la página regular, el filler explícito \(h_*=-vf\), la
autodualidad reducida \(X^{\mathrm{orth}}\), los mapas de cadenas del
triángulo de localización y las bases orientadas de torsión y período. La equivalencia
Kato--FERS y su normalización en uno, el rango completo de \(D_E\), la
exactitud de \eqref{eq:amp-bsd-exacta-sha-tamagawa} y la finitud de
\(\operatorname{Sha}(E)\) son conclusiones de los lemas anteriores, no
premisas independientes. El criterio generador a generador del lema
\ref{lem:amp-bsd-bases-emparejadas-kato-fers} es un control no gobernante
que audita localmente el primer comparador. Los criterios que refutarían una flecha concreta
son:
\begin{enumerate}
 \item un coeficiente raíz distinto de uno, que introduciría una unidad
       relativa entre \(z_K\) y \(z_{\mathrm{PT}}\);
 \item un \(j_q\) no invertible o la omisión de una página Bockstein, que
       alteraría el grado o el determinante de altura;
 \item un divisor de Smith nulo, que impediría la finitud de \(F_E\);
 \item el fallo de exactitud de
       \eqref{eq:amp-bsd-exacta-sha-tamagawa}, que rompería la identificación
       de los factores finitos;
 \item una inversión de orientación en torsión o período, que cambiaría el
       elemento basado aunque preservase su línea no orientada.
\end{enumerate}
La proyección de Manin sin historia y el reescalamiento visible módulo tres
violan el primer criterio; por ello constituyen pruebas de necesidad para
esos modelos reducidos y no premisas del teorema
\ref{thm:amp-bsd-comparadores-construidos}.


\begin{theorem}[Comparación determinantal local--global]
\label{thm:bsd-suc-local-global}
Sean los comparadores basados ya construidos en el
\cref{thm:amp-bsd-comparadores-construidos}:
Kato--FERS/Iwasawa,
Poitou--Tate/Poincaré--Néron--Tate,
Smith--Tamagawa--Shafarevich, torsión y período.
Si \(d_{\mathrm{an}}\) es el grado de la sección analítica en la
especialización central y \(d_{\mathrm{ar}}\) el grado de la línea Selmer
basada, entonces
\begin{equation}\label{eq:bsd-suc-grados}
 d_{\mathrm{an}}=d=d_{\mathrm{ar}},
\end{equation}
y se obtiene la igualdad de elementos distinguidos
\begin{equation}\label{eq:bsd-suc-linea-final}
 \mathfrak z_{\mathrm{an}}(E)=\mathfrak z_{\mathrm{ar}}(E)
 \quad\text{en}\quad\mathcal L_E.
\end{equation}
El miembro analítico es la publicación del término inicial; el miembro
aritmético es la publicación del regulador completo, los factores
elementales de Smith,
los factores locales, la torsión y el período.
\end{theorem}

\begin{proof}
La construcción basada da
\([\Phi_{K/\mathrm{FERS}}^i](T)=I+TB_i(T)\) en cada grado. El comparador
basado de Kato--FERS expresa entonces la sección analítica corregida
como
\begin{equation}\label{eq:bsd-suc-fers}
 L_{\ell}^{\mathrm{corr}}(E,T)
 =\rho_{E,\ell}(T)\det S_E^{\mathrm{corr}}(T),
 \qquad \rho_{E,\ell}(0)=1.
\end{equation}
Por \eqref{eq:bsd-suc-orden}, su grado central es \(d\); por
\eqref{eq:bsd-suc-lt}, su término inicial es exactamente
\(R_{\mathrm{Boc}}^{\mathrm{der}}\), no ese regulador multiplicado por una
unidad indeterminada. Los isomorfismos \(j_q\) de
\eqref{eq:bsd-suc-jq} transportan cada \(\tau_q\) al determinante de la
altura derivada de su página. La página regular transporta el sumando
acíclico que las páginas positivas no ven.

La comparación de Poincaré--Néron--Tate identifica el grado de esa línea de
alturas con \(d_{\mathrm{ar}}\). Los factores elementales de Smith descomponen su parte
finita en los factores de Shafarevich y Tamagawa; el cociente de las redes
integrales aporta el cuadrado de la torsión, y el lector arquimediano aporta
el período. Todos estos mapas actúan sobre la misma línea raíz. Por el
teorema \ref{thm:bsd-suc-raiz-filler}, no queda una unidad relativa entre
las publicaciones Kato y Poitou--Tate. La composición preserva el grado y
el elemento basado, lo que prueba simultáneamente
\eqref{eq:bsd-suc-grados} y \eqref{eq:bsd-suc-linea-final}. La fórmula
escalar correspondiente se reserva para el reconocimiento final.
\end{proof}

\subsection*{Hipótesis y alcance}

La cadena BSD utiliza: el estado enriquecido completo y su coálgebra
Alexander--Whitney; el producto fibrado \eqref{eq:bsd-suc-pullback}; las dos
publicaciones construidas sobre la misma unidad; la línea raíz, su base y su
counidad; el cociente reducido perfecto y saturado, junto con
\(X^{\mathrm{orth}}\); el complejo finita--singular completo; las banderas
Bockstein finitas, sus signos de Koszul y la página regular; y las
comparaciones de altura, Smith, Tamagawa, Shafarevich, torsión y períodos
con una orientación común. La verificación grado a grado de las bases
emparejadas y de los residuales \(\varepsilon_i(b)\) queda registrada como
control local de la realización, no como premisa gobernante del cierre.
Bajo esta estructura de partida explícita,
\eqref{eq:bsd-suc-raiz}, \eqref{eq:bsd-suc-hstar},
\eqref{eq:bsd-suc-lt} y \eqref{eq:bsd-suc-linea-final} son igualdades
exactas.

La conservación de unidad, el recorrido de dos hojas y la publicación
local--global se formalizan mediante la coálgebra de cadenas, el producto
fibrado genealógico, la publicación acoplada, la reducción ortogonal y la
torre Bockstein. Las identidades auxiliares verifican por separado la propiedad de elemento grupal,
la conservación de la historia bajo \(\widehat\Pi_0\), la identidad de la
unidad en ambos canales, la involución reducida, la forma normal
finita--singular y la identidad determinantal de orden arbitrario.

\paragraph{Controles de ablación BSD.}
\noindent\textbf{Estatuto: CONTROL\_NO\_GOBERNANTE.}
Las tres ablaciones siguientes reciben las identidades ya demostradas y
falsifican únicamente modelos reducidos. No aportan premisas al propietario,
no condicionan su cierre y no poseen una arista causal de retorno hacia él.

\begin{ablation}[Proyección de Manin sin historia]
\label{abl:bsd-suc-manin}
Sustituir \(P_E^{\mathrm{gen}}\) por \(P_E^{\mathrm{Man}}\) borra datos que
la publicación acoplada necesita. En efecto, la ruta
\(\gamma=(\mathsf N\mathsf E)^{54}\) puede tener la misma celda y la misma
fase visible que la ruta trivial, pero conserva enrollamiento y memoria
distintos en \(G_{m,n}\). El aumento de Manin las identifica; la segunda
proyección de \eqref{eq:bsd-suc-pullback} no.
\end{ablation}

\begin{proof}
Todo lift basado al complejo de Manin factoriza, hasta homotopía, por la
sección del aumento \(s\varepsilon_{\mathrm{AW}}\). Por ello depende sólo
de la coordenada visible. En cambio, la segunda componente de
\(\widehat\Pi_0(c)\) es el propio \(c\), de modo que conserva el simplex de
historia y distingue el enrollamiento de \(\gamma\). La ablación refuta la
proyección desnuda; no refuta el producto fibrado empleado en el teorema.
\end{proof}

\begin{ablation}[Reescalamiento de la raíz visible sólo módulo tres]
\label{abl:bsd-suc-twist}
En la especialización \(c=1\) de
\eqref{eq:bsd-suc-diferencial-normal}, el ciclo alterado
\begin{equation}\label{eq:bsd-suc-twist4}
 z_K^{(4)}=4r+3e+b
\end{equation}
satisface
\begin{equation}\label{eq:bsd-suc-twist-defecto}
 z_{\mathrm{PT}}-z_K^{(4)}
 =\partial(-f)-3r.
\end{equation}
El defecto \(-3r\) desaparece módulo tres y reaparece módulo nueve. Por
consiguiente, \eqref{eq:bsd-suc-twist4} conserva una sombra parcial, pero no
la unidad basada ni \eqref{eq:bsd-suc-raiz}.
\end{ablation}

\begin{proof}
La igualdad \(\partial f=3e+b\) da inmediatamente
\eqref{eq:bsd-suc-twist-defecto}. Módulo tres, el segundo sumando es nulo;
módulo nueve es una clase raíz no divisible por una frontera del complemento.
Además, su coordenada raíz es cuatro, mientras
\eqref{eq:bsd-suc-normalizacion-raiz} exige uno. El ejemplo verifica la
necesidad del par acoplado y de la orientación; no introduce una excepción
al teorema que conserva ambos datos.
\end{proof}

\begin{ablation}[Pérdida de páginas del regulador]
\label{abl:bsd-suc-bockstein}
La matriz \(\operatorname{diag}(T^2,T)\) muestra que el primer Bockstein no
recupera el acarreo de orden superior. La matriz
\(\operatorname{diag}(2,T^2)\) tiene trivialización positiva
\(\tau_2=1\), pero
\(\tau_{\mathrm{reg}}=2=\operatorname{LT}_{T}(S)\). Por tanto, ni la
primera página ni el producto de páginas positivas sustituyen al regulador
completo \eqref{eq:bsd-suc-rboc}.
\end{ablation}

\begin{proof}
En el primer ejemplo, \(d=3\), \(r_0=2\) y
\(c_I^{\deg}=1\); la página uno no contiene ese último grado de memoria. En
el segundo, el único divisor singular es \(T^2\), pero la dirección regular
aporta el factor dos. Omitirla produciría uno en vez de dos. Ambos cálculos
son casos diagonales de \eqref{eq:bsd-suc-smith} y verifican la necesidad de
todas las piezas de \eqref{eq:bsd-suc-rboc}.
\end{proof}
\section*{Reconocimiento matemático final de Birch--Swinnerton--Dyer}

El desarrollo anterior concluye antes de emplear su formulación histórica
como criterio.

En el dominio declarado por el teorema
\ref{thm:bsd-suc-local-global}, la comparación de Smith identifica el
componente finito de la línea aritmética con la descomposición
Shafarevich--Tamagawa y demuestra la finitud de
\(\operatorname{Sha}(E)\). En consecuencia, su orden en la publicación
escalar siguiente es el cardinal del grupo finito identificado, no un símbolo
formal añadido a posteriori.

La traducción escalar de \eqref{eq:bsd-suc-linea-final} proporciona, con
\(r=\operatorname{rank}E(\mathbb Q)\),
\begin{equation}\label{eq:bsd-suc-rango-final}
 \operatorname{ord}_{s=1}L(E,s)=r
\end{equation}
y
\begin{equation}\label{eq:bsd-suc-formula-final}
 \frac{L^{(r)}(E,1)}{r!\,\Omega_E}
 =\frac{
   \lvert\operatorname{Sha}(E)\rvert\,\operatorname{Reg}(E)\prod_v c_v}
  {\lvert E(\mathbb Q)_{\mathrm{tors}}\rvert^2}.
\end{equation}
Esta es la formulación convencional de Birch--Swinnerton--Dyer. El orden,
el término principal y los factores aritméticos no se definen unos mediante
otros: son publicaciones coordinadas de la igualdad basada
\eqref{eq:bsd-suc-linea-final}, cuya raíz común fue fijada previamente por
\eqref{eq:bsd-suc-raiz}.
